%% mg_ch15_21.tex — Chapters 15–21 of Meaning Games (v2)
%% No preamble; \input from Meaning_Games.tex

%% ================================================================
%%  CHAPTER 15: FORMAL POETICS — CONDENSATION, REPETITION, AND RHYME
%% ================================================================

\chapter{Formal Poetics of Idea Composition}
\label{ch:poetics}

\epigraph{``The limits of my language mean the limits of my world.''}{Ludwig Wittgenstein, \textit{Tractatus}, 5.6}

Poetry is the laboratory of language: it is where the ordinary
mechanisms of meaning---composition, resonance, emergence---are
pushed to extremes and made visible.  In this chapter we develop a
\emph{formal poetics} grounded entirely in the algebraic structures
of Idea Diffusion Theory.  Every poetic device we analyse---condensation,
repetition, rhyme, meter, defamiliarisation---will receive a precise
definition in terms of the resonance function~$\rs$, the emergence
term~$\emerge$, and the composition operation~$\compose$.

Wittgenstein famously held that ``what can be shown cannot be said.''
We partially disagree: what a poem \emph{shows} can at least be
\emph{computed}, once the right algebraic framework is in place.

%% ---------------------------------------------------------------
\section{Poetic Condensation}
\label{sec:condensation}

The first poetic phenomenon we formalise is \emph{condensation}:
the packing of a large amount of resonance into a small expression.

\begin{definition}[Condensation ratio]
\label{def:condensation}
Let $a \in \IS$ be an idea and let $|a|$ denote its \emph{syntactic
length} (the number of atomic ideatic constituents in a chosen
decomposition).  The \textbf{condensation ratio} of~$a$ with respect
to a community resonance function~$\rs$ is
\[
  \gamma(a) \;=\; \frac{w(a)}{|a|}
  \;=\; \frac{\rs(a,a)}{|a|}.
\]
An idea~$a$ is \textbf{poetically condensed} if $\gamma(a) >
\bar{\gamma}$, where $\bar{\gamma}$ is the mean condensation ratio
over the ambient ideatic corpus.
\end{definition}

Condensation captures the intuition that a haiku ``says more'' per
syllable than a legal contract.  The next result shows that
composition can increase condensation super-linearly.

\begin{theorem}[Super-linear condensation]
\label{thm:superlinear-condensation}
Suppose $|a \compose b| \le |a| + |b|$ (composition does not
increase syntactic length beyond concatenation).  If the emergence
term satisfies $\emerge(a,b,a \compose b) > 0$, then
\[
  \gamma(a \compose b)
  \;>\;
  \frac{|a|\,\gamma(a) + |b|\,\gamma(b)}{|a|+|b|}.
\]
That is, composition with positive emergence yields condensation
strictly greater than the weighted average of the parts.
\end{theorem}

\begin{proof}
By the decomposition axiom,
\[
  w(a \compose b)
  = \rs(a \compose b,\, a \compose b)
  = \rs(a, a\compose b) + \rs(b, a\compose b)
    + \emerge(a,b,a\compose b).
\]
Since $\rs(a, a\compose b) \ge \rs(a,a) = w(a)$ by non-decreasing
self-resonance under extension (Axiom~A3 of~\cite{idt}), and
similarly $\rs(b, a\compose b) \ge w(b)$, we obtain
\[
  w(a \compose b)
  \;\ge\; w(a) + w(b) + \emerge(a,b,a\compose b)
  \;>\; w(a) + w(b).
\]
Dividing by $|a \compose b| \le |a|+|b|$ gives
\[
  \gamma(a \compose b)
  > \frac{w(a)+w(b)}{|a|+|b|}
  = \frac{|a|\gamma(a)+|b|\gamma(b)}{|a|+|b|}. \qedhere
\]
\end{proof}

\begin{interpretation}
A poet who composes two images with positive emergence creates a
line whose condensation exceeds any convex combination of the
images' individual condensations.  This is the algebraic content
of Ezra Pound's injunction: ``Use no superfluous word.''  Positive
emergence is the mathematical engine of compression.
\end{interpretation}

\begin{remark}[What the proof really says]
Let us unpack the super-linear condensation theorem in plain language.
The result tells us that whenever two ideas are brought together and
something genuinely new arises from their combination---i.e., the
emergence $\emerge(a,b,a \compose b)$ is strictly positive---then the
resulting composition packs more meaning per unit of expression than
either part did alone.

The proof proceeds in two steps.  First, we invoke the \emph{decomposition
axiom}, which breaks the total weight of the composition into three
contributions: (i)~how much each component resonates with the whole,
and (ii)~the emergence surplus.  This decomposition is analogous to
the way the variance of a sum decomposes into individual variances plus
a covariance term.  Second, we note that each component's resonance
with the whole \emph{exceeds} its self-resonance---an idea resonates
at least as strongly with any context containing it as it does with itself.
This is Axiom~A3, the non-decreasing self-resonance principle, which
is the algebraic analogue of the claim that a word means at least as
much in context as it does in isolation.  Combining these two steps,
the total weight exceeds the sum of parts, and dividing by length gives
the super-linear condensation bound.

Philosophically, this is a theorem about the \emph{efficiency of metaphor}.
When a poet writes ``the sea of faith'' (Matthew Arnold), the emergence
between ``sea'' and ``faith'' ensures that the three-word phrase conveys
more per syllable than a paragraph of literal description.  The theorem
quantifies this poetic intuition: metaphor is provably more efficient
than enumeration.
\end{remark}

\begin{interpretation}[Kolmogorov complexity and condensation]
There is an illuminating connection to \emph{Kolmogorov complexity},
the algorithmic measure of information content.  The Kolmogorov
complexity $K(x)$ of a string~$x$ is the length of the shortest
program that produces~$x$.  Our condensation ratio $\gamma(a) = w(a)/|a|$
is a semantic analogue: it measures resonance per unit of syntactic
length.

The super-linear condensation theorem is then the IDT analogue of the
following fact in algorithmic information theory: there exist strings
whose Kolmogorov complexity exceeds the sum of the complexities of their
parts.  In Kolmogorov's world, this happens because the \emph{interaction}
between the parts carries information.  In IDT, the interaction is
precisely the emergence term.

This connection suggests a deeper programme: define the \emph{ideatic
Kolmogorov complexity} $K_{\rs}(a)$ as the minimum syntactic length of
any decomposition $a = a_1 \compose \cdots \compose a_k$ that achieves
weight $w(a)$.  The super-linear condensation theorem then implies that
$K_{\rs}(a \compose b)$ can be strictly less than $K_{\rs}(a) + K_{\rs}(b)$:
composition \emph{compresses} meaning.  This is an open research
direction connecting IDT to algorithmic information theory
(cf.~Open Problem~\ref{op:source-coding}).
\end{interpretation}

%% ---------------------------------------------------------------
\section{Repetition and the Weight Ratchet}
\label{sec:repetition}

Repetition is the simplest poetic device, yet it has non-trivial
algebraic consequences.

\begin{definition}[Iterated self-composition]
\label{def:iterated-composition}
For $a \in \IS$ define the \textbf{$n$-fold repetition}
$a^{\compose n}$ inductively:
\[
  a^{\compose 1} = a,
  \qquad
  a^{\compose (n+1)} = a^{\compose n} \compose a.
\]
\end{definition}

\begin{theorem}[Weight ratchet]
\label{thm:weight-ratchet}
If $\emerge(a^{\compose n}, a, a^{\compose(n+1)}) \ge 0$ for all
$n \ge 1$, then the sequence $(w(a^{\compose n}))_{n \ge 1}$ is
non-decreasing:
\[
  w(a^{\compose 1})
  \;\le\; w(a^{\compose 2})
  \;\le\; w(a^{\compose 3})
  \;\le\; \cdots
\]
If, moreover, $\emerge(a^{\compose n}, a, a^{\compose(n+1)}) > 0$
for all~$n$, the sequence is strictly increasing.
\end{theorem}

\begin{proof}
By induction.  For $n=1$:
\begin{align*}
  w(a^{\compose 2})
  &= \rs(a^{\compose 2}, a^{\compose 2}) \\
  &= \rs(a,a^{\compose 2}) + \rs(a,a^{\compose 2})
     + \emerge(a,a,a^{\compose 2}) \\
  &\ge 2\,\rs(a,a) + 0 = 2\,w(a) \ge w(a).
\end{align*}
The inductive step is identical, replacing $a^{\compose 2}$ with
$a^{\compose(n+1)}$ and using $w(a^{\compose n}) \le
\rs(a^{\compose n}, a^{\compose(n+1)})$.  Strict inequality follows
when each emergence term is positive.
\end{proof}

\begin{interpretation}
The Weight Ratchet theorem explains why refrains, anaphora, and
liturgical repetition \emph{accumulate} cultural weight: each
repetition adds at least as much resonance as the previous one.
A slogan repeated a hundred times has weight bounded below by~$100\,w(a)$.
\end{interpretation}

\begin{remark}
The Weight Ratchet is a \emph{lower bound}.  In practice, audiences
often exhibit \emph{diminishing marginal emergence}: $\emerge(a^{\compose n},a,\cdot) \to 0$ as $n\to\infty$.  This models the
phenomenon of ``semantic satiation''---a word repeated too many times
loses meaning.  We formalise this decay in Section~\ref{sec:meter}.
\end{remark}

\begin{remark}[Proof in plain language]
The weight ratchet theorem is proved by a simple induction, but its
mathematical content is worth spelling out.  At the base case ($n=1$),
we compose $a$ with itself to get $a^{\compose 2}$.  The decomposition
axiom splits $w(a^{\compose 2})$ into two copies of $\rs(a, a^{\compose 2})$
plus the emergence term.  Each copy is at least $\rs(a,a) = w(a)$
by non-decreasing self-resonance, so $w(a^{\compose 2}) \ge 2w(a) + \emerge$.
When emergence is non-negative, $w(a^{\compose 2}) \ge 2w(a) \ge w(a)$.
The inductive step is identical in structure: replacing the base
with the accumulated composition preserves the inequality.

What makes this mathematically interesting is the \emph{monotonicity}:
once weight is gained, it is never lost.  This is a ``ratchet'' in the
precise sense of thermodynamics---the system has a preferred direction.
In the theory of Markov chains, this corresponds to a supermartingale
property of the weight sequence.

From an information-theoretic perspective, the weight ratchet says
that repeated exposure to an idea \emph{accumulates} mutual information
between the idea and its context.  Each repetition is a new ``sample''
that tightens the listener's estimate of the idea's resonance profile.
The Shannon-theoretic analogue is the law of large numbers applied to
the empirical mutual information: $\hat{I}(X;Y) \to I(X;Y)$ as the
sample size grows.  The weight ratchet is thus a semantic law of
large numbers.
\end{remark}

\begin{interpretation}[Repetition in ritual and liturgy]
The weight ratchet provides a mathematical foundation for the
observation, widespread in the anthropology of religion, that
\emph{ritual repetition creates sacredness}.  Rappaport~\cite{rappaport}
argued that the invariance of liturgical form is constitutive of the
sacred: by repeating the same words in the same order, a community
creates a weight that no single utterance could achieve.  The weight
ratchet theorem makes this precise: $w(a^{\compose n}) \ge n \cdot w(a)$
under non-negative emergence, so $n$~repetitions of a ritual formula
create at least $n$~times the weight of a single performance.  The
``sacred'' is, in this analysis, simply the regime where $n$ is very large.
\end{interpretation}

%% ---------------------------------------------------------------
\section{Rhyme as Resonance}
\label{sec:rhyme}

\begin{definition}[Rhyme coefficient]
\label{def:rhyme}
For ideas $a,b \in \IS$ with $w(a), w(b) > 0$, the
\textbf{rhyme coefficient} is the normalised resonance
\[
  \rho(a,b) \;=\;
  \frac{\rs(a,b)}{\sqrt{w(a)\,w(b)}}.
\]
Two ideas \textbf{rhyme} if $\rho(a,b) > \tau$ for a
community-dependent threshold~$\tau \in (0,1]$.
\end{definition}

The rhyme coefficient is a cosine-similarity measure in the
resonance inner-product space.  The next result shows that
Cauchy--Schwarz bounds it.

\begin{proposition}[Cauchy--Schwarz for rhyme]
\label{prop:cs-rhyme}
For all $a,b \in \IS$ with $w(a),w(b)>0$,
\[
  0 \;\le\; \rho(a,b) \;\le\; 1.
\]
Equality $\rho(a,b)=1$ holds if and only if $a$ and $b$ are
\emph{proportional} in the resonance sense: there exists
$\lambda > 0$ such that $\rs(a,c) = \lambda\,\rs(b,c)$ for all~$c$.
\end{proposition}

\begin{proof}
The lower bound follows from positivity of~$\rs$.  The upper bound
is the Cauchy--Schwarz inequality for the bilinear form~$\rs$,
which is positive semi-definite by axiom.  The equality case is the
standard characterisation of equality in Cauchy--Schwarz.
\end{proof}

\begin{example}[Phonological rhyme]
\label{ex:phono-rhyme}
Let $a = \text{``moon''}$ and $b = \text{``June''}$ in an English
poetic corpus.  Their phonological overlap (shared vowel nucleus and
coda) contributes a high $\rs(a,b)$, while their semantic divergence
keeps $w(a)$ and $w(b)$ individually moderate.  The resulting
$\rho(a,b)$ is high---a strong rhyme.  By contrast, $a =
\text{``moon''}$ and $c = \text{``economy''}$ have low phonological
overlap and divergent semantic fields, so $\rho(a,c) \approx 0$:
no rhyme.
\end{example}

\begin{theorem}[Rhyme enhances condensation]
\label{thm:rhyme-condensation}
Let $a,b \in \IS$ with $\rho(a,b) > 0$.  Then
\[
  \gamma(a \compose b)
  \;\ge\;
  \gamma(a) + \gamma(b)
  + \frac{2\,\rho(a,b)\sqrt{w(a)\,w(b)} + \emerge(a,b,a\compose b)}{|a|+|b|}.
\]
In particular, positive rhyme and positive emergence jointly
amplify condensation.
\end{theorem}

\begin{proof}
Expand $w(a\compose b)$ via the decomposition axiom:
\begin{align*}
  w(a\compose b)
  &= \rs(a,a\compose b) + \rs(b,a\compose b)
     + \emerge(a,b,a\compose b) \\
  &\ge w(a) + w(b)
     + 2\,\rs(a,b) + \emerge(a,b,a\compose b).
\end{align*}
Substituting $\rs(a,b) = \rho(a,b)\sqrt{w(a)w(b)}$ and dividing
by $|a|+|b|$ gives the stated bound.
\end{proof}

\begin{remark}[Why rhyme works: a spectral interpretation]
The rhyme coefficient $\rho(a,b)$ is a cosine similarity in the
resonance inner-product space.  In linear algebra, the cosine
similarity between two vectors measures the alignment of their
``directions,'' independent of their magnitudes.  The rhyme
coefficient does the same for ideas: it measures how well $a$ and
$b$ point in the same ``meaning direction,'' regardless of how
weighty each is individually.

When we say that ``moon'' and ``June'' rhyme, we are saying that
their meaning-direction vectors are nearly parallel in the
resonance inner-product space.  The Cauchy--Schwarz bound
$\rho(a,b) \le 1$ is then the geometric statement that no two
distinct directions can be \emph{more} than parallel.  Equality
$\rho = 1$ means the ideas are semantically proportional---they
differ only in intensity, not in kind.

This spectral perspective connects rhyme to the \emph{eigenstructure}
of the resonance operator.  If we define the resonance operator
$R : \IS \to \IS^*$ by $(Ra)(b) = \rs(a,b)$, then the eigenideas
of $R$---the ideas $e$ satisfying $Re = \lambda e$---are the
``pure tones'' of the meaning space.  Every idea is a superposition
of these eigenideas, and rhyme between $a$ and $b$ measures the
overlap of their spectral decompositions.  Ideas that share dominant
eigencomponents rhyme strongly; ideas with disjoint spectra do not
rhyme at all.  This is the spectral theory of emergence previewed
here and developed fully in Chapter~\ref{ch:information}.
\end{remark}

%% ---------------------------------------------------------------
\section{Meter and Emergence Decay}
\label{sec:meter}

\begin{definition}[Emergence decay function]
\label{def:emergence-decay}
An \textbf{emergence decay function} for an idea~$a$ is a function
$\delta_a : \N \to [0,\infty)$ such that
\[
  \emerge(a^{\compose n}, a, a^{\compose(n+1)})
  \;=\; \delta_a(n) \cdot \emerge(a,a,a^{\compose 2})
\]
with $\delta_a(1) = 1$ and $\delta_a$ non-increasing.
\end{definition}

\begin{theorem}[Bounded weight under exponential decay]
\label{thm:bounded-weight}
If $\delta_a(n) = \alpha^{n-1}$ for some $0 < \alpha < 1$, then
\[
  w(a^{\compose n})
  \;\le\; n\,w(a) + \frac{\emerge_0}{1-\alpha},
\]
where $\emerge_0 = \emerge(a,a,a^{\compose 2})$.  In particular,
$w(a^{\compose n}) = O(n)$ as $n \to \infty$.
\end{theorem}

\begin{proof}
Summing the telescoping contributions:
\begin{align*}
  w(a^{\compose n})
  &\le n\,w(a) + \sum_{k=1}^{n-1} \delta_a(k)\,\emerge_0 \\
  &= n\,w(a) + \emerge_0 \sum_{k=0}^{n-2} \alpha^k \\
  &\le n\,w(a) + \frac{\emerge_0}{1-\alpha}. \qedhere
\end{align*}
\end{proof}

\begin{interpretation}
Meter is emergent rhythm: the ear expects a pattern, and
each additional repetition of the metrical foot adds less
novelty (decreasing $\delta_a$).  Exponential decay models
the rapid habituation that makes iambic pentameter feel
``natural'' after two or three feet.  The bound $O(n)$ says
that weight grows at most linearly with length---a poem
cannot cheat by adding empty repetitions.
\end{interpretation}

\begin{remark}[Proof in plain language: why repeated emergence decays]
The bounded weight theorem is a convergence result disguised as an
inequality.  Here is the intuition.

Imagine a drum being struck repeatedly.  The first strike creates a
sharp, novel sound (high emergence).  The second strike is still novel
but less so (the ear is partially habituated).  By the tenth strike,
the novelty has largely vanished, and each subsequent strike adds
only the base weight of the drum itself.

Mathematically, the exponential decay $\delta_a(n) = \alpha^{n-1}$
models this habituation.  The total accumulated emergence from $n$
repetitions is $\emerge_0 \sum_{k=0}^{n-2} \alpha^k$, a geometric
series that converges to $\emerge_0 / (1-\alpha)$ as $n \to \infty$.
The total weight is therefore bounded by $n \cdot w(a)$ (the linear
contribution of each repetition's base weight) plus a constant
(the finite total emergence).

This $O(n)$ bound is sharp: without emergence decay, the weight
could grow super-linearly (by the weight ratchet theorem).  Decay
constrains growth to linearity.  The constant $\emerge_0/(1-\alpha)$
measures the ``total novelty budget'' of repetition: the amount of
emergence that a repeated idea can ever create, summed over all
future repetitions.

From a rate-distortion perspective (Section~\ref{sec:rate-distortion}),
exponential emergence decay means that the rate at which a repeated
idea generates information \emph{decreases geometrically}.  After
sufficiently many repetitions, the rate drops below the channel
capacity, and further repetitions transmit no additional information.
This is the information-theoretic content of boredom.
\end{remark}

%% ---------------------------------------------------------------
\section{Defamiliarisation as Emergence Shock}
\label{sec:defamiliarisation}

Viktor Shklovsky's notion of \emph{ostranenie} (making strange)
receives a clean formalisation.

\begin{definition}[Defamiliarisation index]
\label{def:defamiliarisation}
Let $a,b \in \IS$ with $\rho(a,b) \approx 0$ (semantically
distant ideas).  The \textbf{defamiliarisation index} of the
composition $a \compose b$ is
\[
  \Delta(a,b)
  \;=\;
  \frac{\emerge(a,b,a\compose b)}{w(a)+w(b)}.
\]
A composition is \textbf{defamiliarising} if $\Delta(a,b) > 1$.
\end{definition}

\begin{theorem}[Defamiliarisation amplifies weight]
\label{thm:defam-weight}
If $\Delta(a,b)>1$ and $\rho(a,b) \ge 0$, then
\[
  w(a \compose b) > 2\bigl(w(a)+w(b)\bigr).
\]
\end{theorem}

\begin{proof}
From the decomposition axiom,
\[
  w(a\compose b)
  \ge w(a) + w(b) + 2\rs(a,b) + \emerge(a,b,a\compose b).
\]
Since $\rs(a,b)\ge 0$ and $\emerge(a,b,a\compose b) =
\Delta(a,b)(w(a)+w(b)) > w(a)+w(b)$, we get
$w(a\compose b) > 2(w(a)+w(b))$.
\end{proof}

\begin{interpretation}
When a poet juxtaposes two unrelated images---the ``Petals on a wet,
black bough'' of Pound's \emph{Metro} poem---the emergence dominates
the resonance.  The composition weighs more than twice its parts.
Defamiliarisation is a \emph{super-additive} shock.
\end{interpretation}

\begin{remark}[Proof explained: why defamiliarisation doubles weight]
The defamiliarisation theorem has a beautifully simple proof that
deserves unpacking.  The key insight is that when $\Delta(a,b) > 1$,
the emergence term alone exceeds the total weight of the parts:
$\emerge(a,b,a \compose b) > w(a) + w(b)$.  Combined with the
decomposition axiom---which tells us that the weight of the whole
\emph{also} includes the resonances $\rs(a, a\compose b)$ and
$\rs(b, a\compose b)$, each at least $w(a)$ and $w(b)$ respectively---we
get $w(a \compose b) > 2(w(a) + w(b))$.

Philosophically, this is a theorem about the \emph{power of surprise}.
When two ideas have no prior resonance ($\rho \approx 0$), their
combination must derive all its weight from emergence---from the
genuinely new.  The condition $\Delta > 1$ says that this novelty
is so powerful that it exceeds the combined weight of the familiar.
This is the algebraic content of Shklovsky's insight: defamiliarisation
does not merely \emph{add} to meaning; it \emph{multiplies} it.

From the perspective of topology (Chapter~\ref{ch:topology}), the
defamiliarisation index $\Delta(a,b)$ is large precisely when $a$
and $b$ lie in distant regions of the resonance metric space:
$d_{\rs}(a,b) \approx \sqrt{2}$.  The theorem then says that
composing topologically distant ideas creates disproportionate weight.
This is a form of \emph{long-range interaction} in the meaning space,
analogous to the way entangled particles in quantum mechanics exhibit
correlations that exceed any local explanation.
\end{remark}

\begin{lstlisting}
-- Lean 4 sketch: Defamiliarisation bound
theorem defamiliarisation_amplifies
    (a b : Idea) (w : Idea -> Real)
    (rs : Idea -> Idea -> Real)
    (emerge : Idea -> Idea -> Idea -> Real)
    (hw : forall x, w x = rs x x)
    (hdecomp : w (comp a b) >= w a + w b + 2 * rs a b + emerge a b (comp a b))
    (hrs_nn : rs a b >= 0)
    (hDelta : emerge a b (comp a b) > w a + w b) :
    w (comp a b) > 2 * (w a + w b) := by
  have h1 := hdecomp
  linarith
\end{lstlisting}

%% ---------------------------------------------------------------
\section{The Poetic Function: A Synthesis}
\label{sec:poetic-function}

We can now unify the devices above into a single functional.

\begin{definition}[Poetic functional]
\label{def:poetic-functional}
For a sequence of ideas $\mathbf{a} = (a_1,\ldots,a_n)$, define
the \textbf{poetic functional}
\[
  \Pi(\mathbf{a})
  \;=\;
  \underbrace{\gamma(a_1\compose\cdots\compose a_n)}_{\text{condensation}}
  \;+\;
  \lambda_\rho \sum_{i<j} \rho(a_i,a_j)
  \;+\;
  \lambda_\Delta \sum_{i} \Delta(a_i,a_{i+1}),
\]
where $\lambda_\rho, \lambda_\Delta > 0$ are weights reflecting
a community's aesthetic preferences.
\end{definition}

\begin{remark}
The poetic functional $\Pi$ is \emph{not} a utility function in
the game-theoretic sense (that comes in Chapter~\ref{ch:bayesian}).
It is a \emph{descriptive} measure: high~$\Pi$ correlates with what
communities recognise as ``good poetry.''  The normative question---how
\emph{should} one compose?---requires the mechanism-design tools
of Chapter~\ref{ch:mechanism}.
\end{remark}

%% ---------------------------------------------------------------
\section{Spectral Decomposition of Poetic Form}
\label{sec:spectral-poetics}

The resonance function~$\rs$ defines a positive semi-definite bilinear
form on~$\IS$.  As in any inner-product space, we may seek an
\emph{eigendecomposition}---a basis of ``pure ideas'' in terms of which
every idea can be expressed.

\begin{definition}[Resonance operator]
\label{def:resonance-operator}
The \textbf{resonance operator} $R : \IS \to \IS^*$ is defined by
$(Ra)(b) = \rs(a,b)$.  An idea $e \in \IS$ is an \textbf{eigenidea}
of $R$ with eigenvalue $\lambda \ge 0$ if $\rs(e,b) = \lambda \cdot
\langle e, b \rangle$ for all $b$ in a suitable sense.  We write
$\mathrm{Spec}(\IS)$ for the spectrum of~$R$.
\end{definition}

\begin{theorem}[Spectral theorem for ideatic spaces]
\label{thm:spectral-poetics}
If $\IS$ is a separable Hilbert space under the resonance inner product,
then the resonance operator $R$ has a countable spectrum $\{\lambda_k\}_{k \ge 1}$
with $\lambda_1 \ge \lambda_2 \ge \cdots \ge 0$ and corresponding
orthonormal eigenideas $\{e_k\}$.  Every idea $a \in \IS$ admits a
spectral decomposition
\[
  a = \sum_{k=1}^{\infty} \alpha_k(a)\, e_k,
  \qquad \alpha_k(a) = \rs(a, e_k)/\lambda_k,
\]
and the weight is $w(a) = \sum_k \lambda_k\, \alpha_k(a)^2$.
\end{theorem}

\begin{proof}
Since $R$ is a positive semi-definite, self-adjoint operator on a
separable Hilbert space, the spectral theorem for compact operators
(or, more generally, the spectral theorem for bounded self-adjoint
operators) yields the stated decomposition.  The weight formula follows
from $w(a) = \rs(a,a) = \sum_k \lambda_k \alpha_k^2$ by Parseval's
identity in the resonance inner product.
\end{proof}

\begin{interpretation}[Eigenideas as archetypal themes]
The eigenideas $\{e_k\}$ are the \emph{archetypes} of a meaning
space---the irreducible thematic elements from which all ideas are
composed.  The eigenvalue $\lambda_k$ measures the ``resonance power''
of the $k$-th archetype: large $\lambda_k$ means the archetype
resonates strongly across the community.

In a literary corpus, we might expect the leading eigenideas to
correspond to universal themes---love, death, conflict, redemption---while
the tail eigenideas capture niche or esoteric concerns.  The spectral
decomposition of a poem $a = \sum \alpha_k e_k$ then reveals its
\emph{thematic fingerprint}: the coefficients $\alpha_k$ tell us how
much of each archetype the poem invokes.

This connects IDT to \emph{latent semantic analysis} (LSA) and to the
singular value decomposition of term-document matrices in computational
linguistics.  The key difference is that IDT's eigendecomposition is
\emph{axiomatically grounded}---it follows from the structure of the
resonance function, not from a statistical model.  As proved in
Volume~I, the resonance axioms guarantee that $R$ is positive
semi-definite, so the spectral theorem applies unconditionally.
\end{interpretation}

\begin{definition}[Spectral condensation]
\label{def:spectral-condensation}
The \textbf{spectral condensation} of an idea~$a$ is the effective
number of eigenideas it activates:
\[
  N_{\mathrm{eff}}(a) = \exp\left(-\sum_k p_k(a) \log p_k(a)\right),
  \qquad p_k(a) = \frac{\lambda_k \alpha_k(a)^2}{w(a)}.
\]
A poem with low $N_{\mathrm{eff}}$ is thematically focused; one with
high $N_{\mathrm{eff}}$ is thematically diffuse.
\end{definition}

\begin{remark}
The spectral condensation $N_{\mathrm{eff}}$ is the exponential of the
Shannon entropy of the spectral distribution $\{p_k\}$.  This connects
the poetic notion of thematic focus to information theory: a focused
poem has \emph{low entropy} in spectral space.  Great poetry, in this
framework, achieves high condensation ratio~$\gamma$ (much meaning per
syllable) with low spectral condensation~$N_{\mathrm{eff}}$ (few
archetypal themes)---it says a great deal about very few things.
\end{remark}



%% ================================================================
%%  CHAPTER 16: BAYESIAN MEANING GAMES
%% ================================================================

\chapter{Bayesian Meaning Games}
\label{ch:bayesian}

\epigraph{``In the great chess-Loss of life, the pawns of the soul
go forward, and may not move back.''}{adapted from Wittgenstein,
\textit{Culture and Value}}

Classical game theory assumes that payoffs are common knowledge.
In the space of ideas, this assumption is absurd: no speaker knows
in advance how an audience will weigh an utterance.  We therefore
pass to \emph{Bayesian} (incomplete-information) games, where
each player holds a private \emph{type} drawn from a prior
distribution over possible resonance functions.

This chapter constructs Bayesian meaning games, proves existence
of Bayesian Nash equilibria in the ideatic setting, and shows that
iterated composition produces Bayesian updating---the audience
\emph{learns} the speaker's type by observing successive compositions.

%% ---------------------------------------------------------------
\section{Types and Priors}
\label{sec:types}

\begin{definition}[Ideatic type]
\label{def:ideatic-type}
An \textbf{ideatic type} of a player~$i$ is a resonance function
$\rs_i : \IS \times \IS \to \R$ satisfying the IDT axioms.
The \textbf{type space} of player~$i$ is a measurable space
$(\Theta_i, \mathcal{F}_i)$ equipped with a prior probability
measure $\mu_i \in \mathcal{P}(\Theta_i)$.
\end{definition}

\begin{definition}[Bayesian meaning game]
\label{def:bayesian-game}
A \textbf{Bayesian meaning game} is a tuple
\[
  G^B = \bigl(N,\; \IS,\; (\Theta_i)_{i\in N},\;
  (\mu_i)_{i\in N},\; (u_i)_{i\in N}\bigr),
\]
where
\begin{itemize}
  \item $N = \{1,\ldots,n\}$ is the player set,
  \item $\IS$ is the shared ideatic space with $\compose$,
  \item $\Theta_i$ is the type space for player~$i$,
  \item $\mu_i \in \mathcal{P}(\Theta_i)$ is the common prior
        on player~$i$'s type,
  \item $u_i : \IS^n \times \Theta_i \to \R$ is
        $u_i(a_1,\ldots,a_n;\theta_i)
        = \rs_{\theta_i}(a_i,\, a_1\compose\cdots\compose a_n)$.
\end{itemize}
A \textbf{strategy} for player~$i$ is a measurable function
$\sigma_i : \Theta_i \to \IS$.
\end{definition}

\begin{theorem}[Existence of Bayesian Nash equilibrium]
\label{thm:bne-existence}
Let $G^B$ be a Bayesian meaning game in which $\IS$ is compact
in the resonance topology (Definition~\ref{def:resonance-topology})
and each $u_i$ is continuous in all arguments.  Then $G^B$ admits
a Bayesian Nash equilibrium in mixed strategies.
\end{theorem}

\begin{proof}
We apply Milgrom and Weber's theorem on distributional strategies
\cite{milgrom-weber}.  The strategy space $\mathcal{P}(\IS)$ is
compact and convex in the weak-$*$ topology (since $\IS$ is compact
metrizable).  The expected payoff
\[
  U_i(\sigma)
  = \int_{\Theta_i} \int_{\IS^n}
    u_i(a_1,\ldots,a_n;\theta_i)\,
    \prod_{j\neq i} d\sigma_j(a_j \mid \theta_j)\,
    d\sigma_i(a_i \mid \theta_i)\, d\mu_i(\theta_i)
\]
is continuous in $\sigma$ by dominated convergence (since $u_i$ is
bounded on the compact domain) and linear in $\sigma_i$.
Kakutani's fixed-point theorem yields a fixed point of the
best-response correspondence, which is a BNE.
\end{proof}

\begin{remark}[What the existence proof tells us philosophically]
The existence of Bayesian Nash equilibrium in meaning games is not
merely a technical result---it is a philosophical claim about the
\emph{possibility of communication}.  The theorem says: whenever ideas
form a compact space (a finite ``universe of discourse'') and resonance
varies continuously, there is always a stable way for speakers to
compose ideas, even when each player's private aesthetic is unknown to
the others.  Communication is always \emph{possible} in principle.

The proof works by a \emph{fixed-point argument}: it defines a
``best-response map'' (given what everyone else is doing, what should I
say?) and shows that this map has a fixed point (a situation where
everyone's response is best, given everyone else's response).  The key
mathematical ingredient is Kakutani's theorem, which guarantees fixed
points for set-valued maps on compact convex sets.  The compactness of
the ideatic space ensures that the space of strategies is well-behaved;
the continuity of resonance ensures that small changes in strategy
produce small changes in payoff.

This echoes a deep theme in the philosophy of language: meaning is not
determined by individual intention alone, but by an equilibrium of
mutual expectations.  Wittgenstein's ``language games'' are precisely
such equilibria.  The existence theorem says that language games
always have at least one solution, though it does not say that the
solution is unique or optimal (see Open Problem~\ref{op:bne-unique}).
\end{remark}

\begin{interpretation}[Connection to Bayesian persuasion]
The Bayesian meaning game framework naturally connects to the
theory of \emph{Bayesian persuasion} developed by Kamenica and
Gentzkow~\cite{kamenica-gentzkow}.  In their framework, a sender
designs an information structure (a ``signal'') to influence a
receiver's action.  In IDT, the ``signal'' is a composed idea
$a \in \IS$, and the ``information structure'' is the mapping from
the speaker's type $\theta$ to the idea $\sigma(\theta)$.

The Kamenica--Gentzkow insight is that the sender's optimal strategy
depends on the \emph{concavification} of the sender's payoff as a
function of the receiver's posterior belief.  In IDT terms: the speaker
chooses ideas to maximise the concave envelope of the resonance
payoff, viewed as a function of the listener's posterior type belief.
This connects persuasion to the \emph{geometry} of the type space:
the sender exploits convexity gaps in the resonance landscape.

As proved in Volume~I, the resonance function $\rs$ is positive
semi-definite, which implies that the payoff function $u_i$ is
concave in certain directions but not all.  The non-concave directions
are precisely where persuasion has bite: the speaker can improve her
payoff by strategically withholding or combining ideas.
\end{interpretation}

%% ---------------------------------------------------------------
\section{Learning Through Iterated Composition}
\label{sec:learning}

\begin{definition}[Composition history]
\label{def:composition-history}
A \textbf{composition history} of length~$T$ is a sequence
$h_T = (a_1, a_2, \ldots, a_T) \in \IS^T$ of ideas composed
sequentially.  The \textbf{cumulative composition} is
$c_T = a_1 \compose a_2 \compose \cdots \compose a_T$.
\end{definition}

\begin{definition}[Posterior type belief]
\label{def:posterior}
After observing history $h_T$ from player~$j$, player~$i$ forms
the posterior
\[
  \mu_j^{(T)}(\theta_j \mid h_T)
  \;\propto\;
  \mu_j(\theta_j)
  \prod_{t=1}^{T}
  \ell(\,a_t \mid \theta_j,\, h_{t-1}\,),
\]
where $\ell(a_t \mid \theta_j, h_{t-1})$ is the likelihood of
player~$j$ with type~$\theta_j$ choosing idea~$a_t$ given
past history~$h_{t-1}$.
\end{definition}

\begin{theorem}[Bayesian consistency of iterated composition]
\label{thm:bayesian-consistency}
If the likelihood model satisfies the identifiability condition
\[
  \theta \neq \theta' \;\Longrightarrow\;
  \exists\, a \in \IS:\;
  \ell(a \mid \theta, h) \neq \ell(a \mid \theta', h)
  \quad\text{for some } h,
\]
then the posterior $\mu_j^{(T)}$ converges weakly to the Dirac
measure $\delta_{\theta_j^*}$ at the true type~$\theta_j^*$,
$\mu_j$-almost surely, as $T \to \infty$.
\end{theorem}

\begin{proof}
This is a consequence of Doob's posterior consistency theorem
applied to the product measure on $\IS^{\N}$ induced by the
likelihood.  Identifiability ensures that the KL divergence
$D_{\mathrm{KL}}(\ell(\cdot\mid\theta^*) \| \ell(\cdot\mid\theta)) > 0$
for $\theta \neq \theta^*$, which is the sufficient condition
for Doob's theorem.
\end{proof}

\begin{interpretation}
Iterated composition is the engine of mutual understanding:
every new utterance reveals something about the speaker's
resonance function.  Over time, the listener converges to the
true type.  This is why long conversations produce deeper
understanding than one-shot exchanges---the Bayesian posterior
tightens with each composition.
\end{interpretation}

\begin{remark}[How the proof works: Doob meets Wittgenstein]
The Bayesian consistency theorem draws on one of the deepest results
in probability theory: Doob's theorem on posterior consistency (1949).
Doob showed that for \emph{almost all} true parameters $\theta^*$
(with respect to the prior $\mu$), the posterior distribution concentrates
on $\theta^*$ as data accumulates.  The ``almost all'' qualifier is
crucial---it means that consistency can fail only on a set of prior
measure zero.

In our setting, each composed idea $a_t$ is a ``data point'' about the
speaker's type.  The identifiability condition ensures that different
types produce distinguishably different compositions---no two resonance
functions generate the same distribution over ideas.  Under this condition,
the Kullback--Leibler divergence between the true likelihood and any
alternative is strictly positive: $D_{\mathrm{KL}}(\ell(\cdot \mid \theta^*)
\| \ell(\cdot \mid \theta)) > 0$ for $\theta \neq \theta^*$.  Doob's
theorem then guarantees that the posterior concentrates.

The philosophical content is striking.  Wittgenstein argued that we
cannot directly access another person's ``form of life''---their private
experience of meaning.  The Bayesian consistency theorem partially
vindicates this concern (we never observe the type directly) while
simultaneously undermining it (we converge to the truth through
observation).  Understanding another person's meaning is not an act of
telepathic insight; it is a process of \emph{statistical inference} that
succeeds under identifiability.  The ``depth'' of understanding that
comes from long acquaintance is, mathematically, the tightness of a
Bayesian posterior.
\end{remark}

\begin{remark}[Prior composition through emergence]
A subtle but important feature of the Bayesian framework is the way
priors \emph{compose} with evidence.  In standard Bayesian inference,
the prior $\mu(\theta)$ and the likelihood $\ell(x \mid \theta)$ combine
via Bayes' rule to produce the posterior $\mu(\theta \mid x) \propto
\mu(\theta) \ell(x \mid \theta)$.  In IDT, the ``evidence'' is a composed
idea, and the likelihood depends on the \emph{emergence} generated by
composition.

Specifically, consider a listener with prior $\mu_j(\theta_j)$ who
observes the speaker compose $a_1$ followed by $a_2$.  The likelihood
of observing $a_2$ after $a_1$ depends on the emergence
$\emerge(a_1, a_2, a_1 \compose a_2)$---the novelty created by the
sequential composition.  High emergence indicates that $a_2$ was
chosen to create surprise, which is informative about the speaker's
type (surprising compositions are more likely from types with high
emergence capacity).  Low emergence indicates routine composition,
which is less informative.

This means that the \emph{rate} of Bayesian learning depends on the
emergence of the composition sequence.  Speakers who compose with
high emergence are learned faster; speakers who compose routinely
are learned slowly.  This has an elegant cultural interpretation:
\emph{creative speakers are better understood than conventional ones},
because their compositions are more statistically informative.
\end{remark}

%% ---------------------------------------------------------------
\section{Bayesian Persuasion in Ideatic Spaces}
\label{sec:persuasion}

The Kamenica--Gentzkow framework of Bayesian persuasion has a
natural IDT incarnation.

\begin{definition}[Ideatic persuasion problem]
\label{def:persuasion}
An \textbf{ideatic persuasion problem} consists of:
\begin{itemize}
  \item A \emph{sender} (speaker) who observes a state of the world
        $\omega \in \Omega$ and chooses an idea $a = \sigma(\omega) \in \IS$;
  \item A \emph{receiver} (listener) who observes $a$ but not $\omega$,
        forms a posterior $\mu(\omega \mid a)$, and takes an action
        $b = b^*(\mu) \in \IS$ maximising $\rs_2(b, a \compose b)$;
  \item The sender's payoff is $v(\omega, b) = \rs_{\omega}(a, a \compose b)$.
\end{itemize}
The sender's problem is to choose the \emph{information policy}
$\sigma : \Omega \to \IS$ that maximises $\mathbb{E}_\omega[v(\omega,
b^*(\mu(\cdot \mid \sigma(\omega))))]$.
\end{definition}

\begin{theorem}[Concavification principle for ideatic persuasion]
\label{thm:concavification}
The sender's optimal payoff in the ideatic persuasion problem equals
$\mathrm{cav}(V)(\mu_0)$, where $\mu_0$ is the prior belief over
$\Omega$, $V(\mu) = \mathbb{E}_\mu[v(\omega, b^*(\mu))]$ is the
sender's indirect payoff, and $\mathrm{cav}(V)$ denotes the
\emph{concave closure} (the pointwise smallest concave function
dominating $V$).
\end{theorem}

\begin{proof}
This follows from the general Kamenica--Gentzkow concavification result.
Any information policy $\sigma$ induces a distribution over posteriors
$\mu$ that is a mean-preserving spread of the prior $\mu_0$ (by
Bayes-plausibility: $\mathbb{E}[\mu(\cdot \mid a)] = \mu_0$).
The sender maximises $\mathbb{E}[V(\mu)]$ subject to this constraint.
By the supporting hyperplane theorem for concave functions, the
maximum is $\mathrm{cav}(V)(\mu_0)$.
\end{proof}

\begin{interpretation}
The concavification principle tells us that persuasion is a problem
of \emph{exploiting non-concavities} in the payoff landscape.  When
the sender's indirect payoff $V(\mu)$ is already concave in the
listener's belief, persuasion has no value: the sender cannot do better
than full transparency.  But when $V$ has convex regions---when the
sender benefits from the listener's uncertainty---then strategic
composition of ideas can strictly improve the sender's payoff.

In cultural terms, this explains the role of \emph{ambiguity} in
persuasion.  A politician who speaks ambiguously is not merely being
evasive; she is optimally exploiting non-concavities in the payoff
landscape.  The concavification theorem gives her the precise recipe:
choose ideas whose composition creates a distribution over listener
beliefs that maximises her expected resonance.
\end{interpretation}

%% ---------------------------------------------------------------
\section{Signalling Equilibria}
\label{sec:signalling}

\begin{definition}[Signalling meaning game]
\label{def:signalling-game}
A \textbf{signalling meaning game} is a two-player Bayesian
meaning game where player~1 (the \emph{speaker}) has private
type $\theta_1 \in \Theta_1$ and player~2 (the \emph{listener})
has a known type.  The speaker moves first by choosing
$a \in \IS$; the listener observes~$a$ and chooses a response
$b \in \IS$.  Payoffs are
\begin{align*}
  u_1(a,b;\theta_1)
  &= \rs_{\theta_1}(a, a \compose b), \\
  u_2(a,b)
  &= \rs_2(b, a \compose b).
\end{align*}
\end{definition}

\begin{theorem}[Separating equilibrium under monotone resonance]
\label{thm:separating}
Suppose the type space $\Theta_1 = \{\theta_L, \theta_H\}$ with
$w_{\theta_H}(a) > w_{\theta_L}(a)$ for all $a \neq \void$.
If the speaker's strategy space is rich enough (for every type,
the best response is unique), then the signalling game admits a
\textbf{separating equilibrium} in which each type sends a
distinct idea.
\end{theorem}

\begin{proof}
The single-crossing property holds: if $\theta_H$ values
composition more than $\theta_L$, then for any two ideas
$a' \succ a$ (in the weight ordering), the incremental payoff
$u_1(a',b;\theta_H) - u_1(a,b;\theta_H)
> u_1(a',b;\theta_L) - u_1(a,b;\theta_L)$.
By the Spence--Mirrlees condition, the best-response
correspondence for $\theta_H$ lies strictly above that for
$\theta_L$, ensuring separation.
\end{proof}

\begin{interpretation}
The separating equilibrium is the algebraic content of
``actions speak louder than words.''  A speaker whose true
resonance function assigns high weight to an idea can credibly
signal this by composing it, because a low-type speaker would
find the cost (in terms of payoff sacrifice) prohibitive.
\end{interpretation}

\begin{remark}[The proof via Spence--Mirrlees: why separation works]
The separating equilibrium theorem relies on the \emph{single-crossing
property}, a condition named after Michael Spence's model of
educational signalling (for which he shared the 2001 Nobel Prize).

The key insight is that high-type and low-type speakers experience
different marginal costs of ``upgrading'' their composition.  For
the high-type speaker (whose resonance function assigns high weight
to all non-void ideas), the marginal benefit of composing a weightier
idea is large.  For the low-type speaker, the same upgrade yields a
smaller benefit.  The single-crossing property says that these
marginal benefit curves cross \emph{exactly once}: at some threshold
idea~$a^*$, the high type is willing to compose while the low type
is not.

In the proof, we formalise this as follows.  For any two ideas
$a' \succ a$ (where $a'$ is weightier), the incremental payoff for
the high type exceeds that for the low type:
\[
  u_1(a', b; \theta_H) - u_1(a, b; \theta_H) > u_1(a', b; \theta_L) - u_1(a, b; \theta_L).
\]
This follows from the assumption that $w_{\theta_H}(a) > w_{\theta_L}(a)$
for all $a \neq \void$, which means the high type's resonance function
is uniformly stronger.

The Spence--Mirrlees condition then guarantees that the best-response
correspondence for $\theta_H$ lies strictly above that for $\theta_L$.
In equilibrium, the listener can infer the speaker's type from the
choice of idea: a high choice reveals a high type, a low choice
reveals a low type.  This is separation.

The cultural interpretation is that costly signalling---composing
ideas that require genuine resonance capacity---is a reliable indicator
of type.  A community that allows costly signalling (e.g., through
peer-reviewed publication, which requires sustained intellectual
effort) can achieve separation.  A community that does not (e.g.,
social media, where the cost of posting is negligible) is stuck
in pooling equilibria or babbling.
\end{remark}

\begin{lstlisting}
-- Lean 4 sketch: Separating equilibrium existence
theorem separating_equilibrium_exists
    (thetaL thetaH : Type_)
    (w_L w_H : Idea -> Real)
    (h_order : forall a, a != void -> w_H a > w_L a)
    (IS : IdeaSpace)
    (h_rich : forall theta, exists! a, is_best_response theta a) :
    exists sigma : Type_ -> Idea,
      sigma thetaL != sigma thetaH
      /\ is_BNE sigma := by
  sorry -- Full proof uses Spence-Mirrlees single-crossing
\end{lstlisting}

%% ---------------------------------------------------------------
\section{Pooling and Babbling}
\label{sec:pooling}

Not all equilibria are separating.

\begin{definition}[Pooling equilibrium]
\label{def:pooling}
A \textbf{pooling equilibrium} in a signalling meaning game is
a strategy profile where all types of the speaker send the same
idea: $\sigma_1(\theta) = a^*$ for all $\theta \in \Theta_1$.
\end{definition}

\begin{proposition}[The void is always a pooling equilibrium]
\label{prop:void-pooling}
If $\rs(\void, b) = 0$ for all $b \in \IS$ (the void resonates
with nothing), then $\sigma_1(\theta) = \void$ for all~$\theta$
is a pooling equilibrium---the \textbf{babbling equilibrium}.
\end{proposition}

\begin{proof}
If the speaker sends $\void$, the listener learns nothing and
plays a best response to the prior.  Deviating to any $a \neq \void$
yields $u_1(a,b;\theta) = \rs_\theta(a, a\compose b)$, but since
the listener's response~$b$ is fixed at the prior-optimal, the
deviation may not improve payoff.  Hence $\void$ is a best
response for all types if the listener's prior-optimal response
dominates any response to an off-path signal.
\end{proof}

\begin{interpretation}
The babbling equilibrium is silence: when no type has an incentive
to break the ice, communication collapses.  This models the
``wall of silence'' in communities where speaking carries
no differential resonance benefit.
\end{interpretation}

%% ---------------------------------------------------------------
\section{The Value of Information in Meaning Games}
\label{sec:value-info}

How much is it worth to learn another player's type before
composing?

\begin{definition}[Value of information]
\label{def:voi}
The \textbf{value of information} for player~$i$ in a Bayesian
meaning game is the difference between the expected payoff under
complete information and under the prior:
\[
  \mathrm{VoI}_i
  = \mathbb{E}_\theta\bigl[u_i\bigl(\sigma^*_{\mathrm{CI}}(\theta);\theta_i\bigr)\bigr]
  - \mathbb{E}_\theta\bigl[u_i\bigl(\sigma^*_{\mathrm{BNE}}(\theta);\theta_i\bigr)\bigr],
\]
where $\sigma^*_{\mathrm{CI}}$ is the complete-information equilibrium
and $\sigma^*_{\mathrm{BNE}}$ is the Bayesian Nash equilibrium.
\end{definition}

\begin{theorem}[Value of information bound]
\label{thm:voi-bound}
In a Bayesian meaning game with type entropy $H(\Theta_j)$, the
value of information satisfies
\[
  \mathrm{VoI}_i
  \le \sqrt{2\, W_{\max}\, H(\Theta_j)},
\]
where $W_{\max} = \sup_{a \in \IS} w(a)$ is the maximum weight.
\end{theorem}

\begin{proof}
By Pinsker's inequality, the total variation distance between the
prior and the Dirac posterior satisfies
$\| \mu_j - \delta_{\theta_j^*}\|_{\mathrm{TV}} \le \sqrt{H(\Theta_j)/2}$.
The payoff difference is bounded by the Lipschitz constant of the
payoff function (at most $W_{\max}$) times the total variation
distance.  Combining gives $\mathrm{VoI}_i \le W_{\max}
\sqrt{H(\Theta_j)/2} \le \sqrt{2 W_{\max} H(\Theta_j)}$.
\end{proof}

\begin{interpretation}
The value of information is bounded by the geometric mean of the
maximum weight and the type entropy.  When types are highly uncertain
(high $H$), learning is very valuable; when the meaning space is
sparse (low $W_{\max}$), learning is less valuable because even
optimal composition creates limited weight.  This connects the
information-theoretic structure of the type space to the algebraic
structure of the ideatic space in a quantitative way.

The bound also implies that in communities with very low type entropy---where
everyone agrees on resonance---information has little value, and
Bayesian games reduce to complete-information games.  This is the
mathematical content of Wittgenstein's observation that in a
``form of life'' where everyone agrees, there is nothing to argue about.
\end{interpretation}



%% ================================================================
%%  CHAPTER 17: MECHANISM DESIGN FOR MEANING
%% ================================================================

\chapter{Mechanism Design for Meaning}
\label{ch:mechanism}

\epigraph{``The philosopher is not a citizen of any community of
ideas.  That is what makes him a philosopher.''}{Wittgenstein,
\textit{Zettel}, \S455}

Wittgenstein's later philosophy insists that meaning is a social
practice: the rules of a language-game are constituted by the
community that plays them.  But this raises a design question:
can a community \emph{choose} its composition rules so that
truth-telling (or, in our framework, sincere resonance reporting)
is incentive-compatible?  This is the province of \emph{mechanism
design}, which we now adapt to ideatic spaces.

%% ---------------------------------------------------------------
\section{The Revelation Problem}
\label{sec:revelation}

\begin{definition}[Ideatic social choice function]
\label{def:scf}
An \textbf{ideatic social choice function} is a mapping
$f : \prod_{i \in N} \Theta_i \to \IS$ that selects a
``communal idea'' given the true type profile.
\end{definition}

\begin{definition}[Mechanism]
\label{def:mechanism}
An \textbf{ideatic mechanism} is a tuple
$\mathcal{M} = (M_1,\ldots,M_n,\; g)$
where $M_i$ is the \textbf{message space} for player~$i$ and
$g : \prod_i M_i \to \IS$ is the \textbf{outcome function}.
A mechanism \textbf{implements} $f$ if there exists an equilibrium
in which $g(\sigma(\theta)) = f(\theta)$ for all type profiles~$\theta$.
\end{definition}

\begin{theorem}[Revelation principle for meaning games]
\label{thm:revelation}
If a social choice function $f$ is implementable by \emph{any}
mechanism in Bayesian Nash equilibrium, then it is implementable
by a \textbf{direct mechanism} in which $M_i = \Theta_i$ and
truth-telling ($\sigma_i(\theta_i) = \theta_i$) is an equilibrium.
\end{theorem}

\begin{proof}
Standard revelation-principle argument.  Given a mechanism
$\mathcal{M}=(M_i,g)$ and an equilibrium $\sigma^*$, define
the direct mechanism $g'(\theta) = g(\sigma^*(\theta))$.
Truth-telling in the direct mechanism yields the same outcome
as $\sigma^*$ in the original mechanism.  Any profitable
deviation from truth-telling in $g'$ would correspond to a
profitable deviation from $\sigma^*$ in $\mathcal{M}$,
contradicting the equilibrium assumption.
\end{proof}

\begin{remark}[The revelation principle in plain language]
The revelation principle is one of the most powerful results in
economic theory, and its application to meaning games deserves
careful explanation.

The theorem says: if you can implement a social choice function
(a rule for choosing communal ideas) using \emph{any} mechanism---however
complex---then you can also implement it using a very simple mechanism
in which every player just honestly reports their type.  The ``direct''
mechanism asks each player: ``What is your resonance function?''  If
truth-telling is an equilibrium of the direct mechanism, then we lose
nothing by restricting attention to direct mechanisms.

The proof is constructive and elegant.  Given any mechanism $\mathcal{M}$
with an equilibrium $\sigma^*$, we build the direct mechanism $g'$ that
\emph{simulates} $\mathcal{M}$: when player~$i$ reports type $\theta_i$,
the mechanism computes $\sigma^*_i(\theta_i)$ (what player~$i$ would
have done in the original mechanism) and then applies the outcome
function~$g$.  Truth-telling in $g'$ is equivalent to playing $\sigma^*$
in $\mathcal{M}$, so truth-telling is indeed an equilibrium.

For IDT, this is philosophically significant.  It says that the
question ``Can a community design institutions to elicit sincere
reports of aesthetic judgment?'' reduces to the question ``Can a
community design a \emph{direct} institution where honesty is
incentive-compatible?''  This is the Hurwicz programme---named after
Leonid Hurwicz, who shared the 2007 Nobel Prize for mechanism design---applied
to the domain of meaning.

Hurwicz's original insight was that economic institutions should be
evaluated not by their stated goals but by the \emph{incentives} they
create.  The revelation principle is the formal tool that reduces
institutional design to incentive analysis.  In IDT, this means
that cultural institutions---journals, academies, peer review---can
be analysed purely in terms of the incentives they create for
sincere resonance reporting.
\end{remark}

%% ---------------------------------------------------------------
\section{Incentive Compatibility}
\label{sec:ic}

\begin{definition}[Incentive-compatible ideatic mechanism]
\label{def:ic}
A direct mechanism $g : \prod_i \Theta_i \to \IS$ is
\textbf{incentive-compatible (IC)} if for every player~$i$
and every pair of types $\theta_i, \theta_i' \in \Theta_i$,
\[
  \mathbb{E}_{\theta_{-i}}\bigl[
    u_i\bigl(g(\theta_i,\theta_{-i});\, \theta_i\bigr)
  \bigr]
  \;\ge\;
  \mathbb{E}_{\theta_{-i}}\bigl[
    u_i\bigl(g(\theta_i',\theta_{-i});\, \theta_i\bigr)
  \bigr].
\]
\end{definition}

\begin{definition}[Transfer payment]
\label{def:transfer}
A \textbf{resonance transfer} is a function
$t_i : \prod_j \Theta_j \to \R$ specifying a payment
(in ``resonance units'') to player~$i$.  The augmented payoff is
\[
  v_i(\theta) = u_i(g(\theta);\theta_i) + t_i(\theta).
\]
\end{definition}

\begin{theorem}[VCG mechanism for ideatic allocation]
\label{thm:vcg}
Define the VCG transfers
\[
  t_i(\theta)
  = \sum_{j \neq i}
    u_j\bigl(g(\theta);\theta_j\bigr)
  - h_i(\theta_{-i}),
\]
where $h_i$ is an arbitrary function of $\theta_{-i}$ and
$g(\theta) = \arg\max_{a \in \IS} \sum_{j} u_j(a;\theta_j)$.
Then the direct mechanism $(g, t_1,\ldots,t_n)$ is
incentive-compatible and the outcome $g(\theta)$ maximises
social welfare.
\end{theorem}

\begin{proof}
Player~$i$'s augmented payoff under truth-telling is
\begin{align*}
  v_i(\theta_i,\theta_{-i})
  &= u_i(g(\theta);\theta_i) + \sum_{j\neq i} u_j(g(\theta);\theta_j) - h_i(\theta_{-i}) \\
  &= \sum_j u_j(g(\theta);\theta_j) - h_i(\theta_{-i}).
\end{align*}
Since $g(\theta)$ maximises $\sum_j u_j(a;\theta_j)$ over $a$,
truth-telling maximises $v_i$---deviating to $\theta_i'$ yields
an outcome $g(\theta_i',\theta_{-i})$ that attains a weakly
lower social welfare, hence a weakly lower $v_i$.
\end{proof}

\begin{interpretation}
The VCG mechanism for ideatic allocation says: if a community can
commit to a rule that selects the idea maximising total resonance,
and charges each member for the ``externality'' their report imposes,
then every member reports sincerely.  This is a precise rebuttal to
the Wittgensteinian worry that private meaning is inaccessible:
under VCG, agents \emph{voluntarily reveal} their resonance
functions.
\end{interpretation}

\begin{remark}[How the VCG proof works: aligning private and social incentives]
The VCG mechanism is named after Vickrey (1961), Clarke (1971), and
Groves (1973).  Its proof is one of the most elegant in mechanism design,
and it works by a trick that deserves exposition.

The key insight is to \emph{align} each player's private incentive
with the social objective.  Under the VCG transfer
$t_i(\theta) = \sum_{j \neq i} u_j(g(\theta); \theta_j) - h_i(\theta_{-i})$,
player~$i$'s augmented payoff becomes
\[
  v_i = u_i(g(\theta); \theta_i) + \sum_{j \neq i} u_j(g(\theta); \theta_j) - h_i(\theta_{-i})
  = \underbrace{\sum_j u_j(g(\theta); \theta_j)}_{\text{social welfare}} - h_i(\theta_{-i}).
\]
The term $h_i(\theta_{-i})$ depends only on others' types, so player~$i$
cannot affect it.  Therefore, player~$i$'s goal reduces to maximising
social welfare---exactly the designer's goal.  Truth-telling achieves
this maximum (since $g(\theta) = \arg\max_a \sum_j u_j(a; \theta_j)$),
so truth-telling is a dominant strategy.

The philosophical content is remarkable.  The VCG mechanism transforms
\emph{selfish} agents into \emph{altruistic} ones: by construction,
each agent maximises her own payoff by maximising everyone's payoff.
The mechanism does not require agents to \emph{be} altruistic; it
creates a situation where selfish behaviour \emph{is} altruistic
behaviour.  This is the ``invisible hand'' made rigorous.

For meaning games specifically, VCG says that a community can elicit
sincere aesthetic judgments by (i)~selecting the idea that maximises
total resonance, and (ii)~charging each member an ``externality tax''
equal to the difference between the social welfare the community
\emph{would} have achieved without that member's report and the welfare
it actually achieves.  The tax internalises the externality: a member
whose report changes the outcome pays for the disruption she causes.

As proved in Volume~I, the VCG mechanism is the \emph{unique}
mechanism (up to the choice of~$h_i$) that achieves both efficiency
and incentive compatibility when the type space is sufficiently rich.
This uniqueness result, due to Green and Laffont~\cite{green-laffont},
strengthens the case for VCG as the canonical institution for
ideatic allocation.
\end{remark}

%% ---------------------------------------------------------------
\section{Individual Rationality and Participation}
\label{sec:ir}

\begin{definition}[Individually rational mechanism]
\label{def:ir}
A mechanism $(g, t)$ is \textbf{individually rational (IR)} if
every player's expected augmented payoff is non-negative under
truth-telling:
\[
  \mathbb{E}_{\theta_{-i}}\bigl[v_i(\theta_i,\theta_{-i})\bigr]
  \;\ge\; 0
  \quad\text{for all } i, \theta_i.
\]
\end{definition}

\begin{theorem}[Myerson--Satterthwaite impossibility for ideas]
\label{thm:myerson-impossibility}
Consider a bilateral meaning game ($n=2$) where player~1 is a
``seller'' of idea~$a$ and player~2 is a ``buyer.''  If the type
distributions $\mu_1, \mu_2$ have overlapping support and the
social choice function $f$ is ex-post efficient, then no
mechanism implementing~$f$ is simultaneously IC, IR, and
budget-balanced ($t_1 + t_2 = 0$).
\end{theorem}

\begin{proof}
Adapt Myerson and Satterthwaite~\cite{myerson-satterthwaite}.
Efficiency requires trade whenever $\rs_{\theta_2}(a,\cdot) >
\rs_{\theta_1}(a,\cdot)$.  IC pins down expected transfers via
the revenue equivalence lemma.  Budget balance and IR jointly
require the expected gains from trade to be non-negative, which
fails when the type supports overlap---specifically, there exist
$\theta_1, \theta_2$ in the overlap region where the seller's
valuation exceeds the buyer's, making efficient trade impossible
without an outside subsidy.
\end{proof}

\begin{interpretation}
Even in the world of ideas, efficiency, incentive compatibility,
and voluntary participation cannot all be achieved simultaneously.
This impossibility result explains why cultural institutions
(journals, galleries, universities) exist: they act as
``subsidisers'' who absorb the budget deficit, enabling trades
of ideas that the free market cannot support.
\end{interpretation}

\begin{remark}[Why the impossibility matters]
The Myerson--Satterthwaite theorem for ideas is perhaps the most
sobering result in this volume.  Let us unpack why.

The theorem establishes a three-way conflict among desirable
properties: (1)~\emph{efficiency}---the community should select ideas
that maximise total resonance; (2)~\emph{incentive compatibility}---each
member should have no reason to misrepresent their aesthetic preferences;
(3)~\emph{individual rationality}---no one should be forced to participate.

The proof shows that when the ``buyer'' and ``seller'' of an idea have
overlapping valuations (there is genuine uncertainty about whether the
trade should occur), no budget-balanced mechanism can simultaneously
satisfy all three properties.  The revenue equivalence lemma pins down
the expected transfers under IC, and these transfers are incompatible
with both IR and budget balance in the overlap region.

Culturally, this means that every institution for idea exchange involves
a \emph{compromise}.  Academic peer review achieves IC (anonymity
reduces strategic manipulation) but sacrifices efficiency (rejection
of good papers) and IR (unpaid labour).  Art markets achieve
participation (voluntary exchange) but sacrifice IC (strategic pricing)
and efficiency (taste bubbles).  No institutional design can escape
the impossibility; the best we can do is choose which property to
sacrifice.

This result also connects to Maskin's work on implementation theory,
which we develop in the next section.  Maskin showed that a social
choice function is implementable in Nash equilibrium if and only if it
satisfies a condition called \emph{monotonicity}.  The Myerson--Satterthwaite
impossibility can be understood as a failure of Maskin monotonicity in
the bilateral setting.
\end{remark}

%% ---------------------------------------------------------------
\section{Maskin Monotonicity and Implementation}
\label{sec:maskin}

Eric Maskin's contribution to mechanism design theory centres on
the concept of \emph{monotonicity}: a necessary condition for a
social choice function to be implementable in Nash equilibrium.

\begin{definition}[Maskin monotonicity for ideatic spaces]
\label{def:maskin-mono}
An ideatic social choice function $f : \prod_i \Theta_i \to \IS$
satisfies \textbf{Maskin monotonicity} if, whenever $f(\theta) = a$
and for all players~$i$ and ideas~$b$:
\[
  u_i(a; \theta_i) \ge u_i(b; \theta_i)
  \;\Longrightarrow\;
  u_i(a; \theta_i') \ge u_i(b; \theta_i'),
\]
then $f(\theta') = a$ as well (where $\theta'$ agrees with $\theta$
on the lower contour sets of~$a$).
\end{definition}

\begin{theorem}[Maskin's theorem for meaning games]
\label{thm:maskin}
If an ideatic social choice function~$f$ is Nash-implementable by
some mechanism, then $f$ satisfies Maskin monotonicity.  Conversely,
if $f$ satisfies Maskin monotonicity and the ``no-veto-power''
condition (no single player can block a unanimously preferred outcome),
then $f$ is Nash-implementable when $|N| \ge 3$.
\end{theorem}

\begin{proof}
\emph{Necessity.}  Suppose $f(\theta) = a$ and $\sigma^*$ is a Nash
equilibrium implementing~$f$.  If the lower contour set of~$a$ weakly
expands from $\theta$ to~$\theta'$ for all players, then $\sigma^*$
remains an equilibrium under~$\theta'$ (no player gains from deviating,
since the set of outcomes they prefer to~$a$ has only shrunk).  Hence
$f(\theta') = g(\sigma^*(\theta')) = a$.

\emph{Sufficiency.}  The canonical construction uses a mechanism where
each player reports a type, an alternative, and an integer.  If all
reports agree, the outcome is~$f(\theta)$.  If exactly one deviates,
monotonicity ensures the outcome is unchanged unless the deviant
can produce a beneficial deviation.  Integer ties are broken by a
modular arithmetic rule.  The no-veto-power condition handles the
remaining edge cases.
\end{proof}

\begin{interpretation}
Maskin monotonicity has a natural cultural interpretation: if the
community selects idea~$a$ when preferences are $\theta$, and then
preferences shift so that $a$ is liked \emph{even more} (relative
to all alternatives), then $a$ should still be selected.  This is
a basic rationality condition for cultural institutions: don't
reject a classic when it becomes even more relevant.

The sufficiency result says that for communities with at least three
members, any monotonic social choice function can be implemented---there
exists an institution that makes it happen.  This is a deep optimism
result: in communities of sufficient size, good cultural outcomes are
attainable by design.

As proved in Volume~I, the resonance axioms endow the preference
ordering with a rich geometric structure that makes Maskin monotonicity
easier to verify than in the general mechanism-design setting.
\end{interpretation}

%% ---------------------------------------------------------------
\section{Optimal Auction for Resonance Slots}
\label{sec:auction}

\begin{definition}[Resonance slot auction]
\label{def:slot-auction}
A \textbf{resonance slot} is a position in a communal composition
sequence.  A \textbf{resonance slot auction} allocates $k$ slots
to $n > k$ competing ideas based on reported types.
\end{definition}

\begin{theorem}[Revenue-optimal auction]
\label{thm:optimal-auction}
The revenue-optimal auction for resonance slots is a
Myerson-type auction~\cite{myerson-auction} where each idea~$a_i$
with reported type~$\theta_i$ receives a \textbf{virtual
resonance} score
\[
  \psi_i(\theta_i)
  = w_{\theta_i}(a_i)
  - \frac{1-F_i(\theta_i)}{f_i(\theta_i)}
    \cdot w_{\theta_i}'(a_i),
\]
and the $k$~ideas with the highest virtual resonance are
allocated slots, paying the threshold virtual resonance.
\end{theorem}

\begin{proof}[Proof sketch]
Apply Myerson's optimal auction theory.  The virtual valuation
$\psi_i$ arises from the usual integration-by-parts trick in the
IC constraint.  Under the regularity condition ($\psi_i$
non-decreasing in $\theta_i$), the allocation rule that
maximises expected virtual surplus is to select the top-$k$ ideas.
\end{proof}

\begin{interpretation}[Why virtual resonance matters]
The revenue-optimal auction introduces the concept of \emph{virtual
resonance}---a modified version of actual resonance that accounts for
the ``information rent'' each agent extracts.  The virtual resonance
$\psi_i(\theta_i) = w_{\theta_i}(a_i) - \frac{1-F_i(\theta_i)}{f_i(\theta_i)}
\cdot w_{\theta_i}'(a_i)$ differs from actual resonance by the
\emph{hazard rate adjustment} $(1-F)/f$.

The hazard rate measures how ``surprising'' a type is: if $\theta_i$ is
in the tail of the distribution (low $f_i$, high $1-F_i$), the hazard
rate is large, and the virtual resonance is substantially less than the
actual resonance.  Rare types pay a large ``information rent'' because
the mechanism designer must leave them enough surplus to prevent them
from imitating more common types.

In cultural terms, the optimal auction for resonance slots is a model
of \emph{editorial selection}.  A journal editor allocating limited
publication slots must balance actual quality (actual resonance) against
strategic reporting (authors know their own work better than the editor).
The virtual resonance formula tells the editor how much to ``discount''
each submission based on the author's position in the quality distribution.
The regularity condition ($\psi_i$ non-decreasing) ensures that the
ranking by virtual resonance agrees with the ranking by actual resonance
for the ``typical'' case; when regularity fails, the editor must use
``ironing'' (pooling nearby types) to restore monotonicity.

As proved in Volume~I, Myerson's theory extends to multi-dimensional
type spaces with substantial complications.  The one-dimensional case
presented here captures the essential economic logic while remaining
analytically tractable.
\end{interpretation}

\begin{lstlisting}
-- Lean 4 sketch: VCG mechanism incentive compatibility
theorem vcg_ic (N : Finset Player) (IS : IdeaSpace)
    (u : Player -> Idea -> Type_ -> Real)
    (g : (Player -> Type_) -> Idea)
    (t : Player -> (Player -> Type_) -> Real)
    (h_welfare : forall theta, g theta = argmax (fun a => Finset.sum N (fun j => u j a (theta j))))
    (h_vcg : forall i theta,
      t i theta = Finset.sum (N.erase i) (fun j => u j (g theta) (theta j)) - h_i i (fun j => theta j)) :
    forall i theta_i theta_i',
      E_neg_i (fun th => u i (g (update th i theta_i)) theta_i + t i (update th i theta_i))
      >= E_neg_i (fun th => u i (g (update th i theta_i')) theta_i + t i (update th i theta_i')) := by
  sorry -- Standard VCG incentive-compatibility proof
\end{lstlisting}



%% ================================================================
%%  CHAPTER 18: INFORMATION AND EMERGENCE
%% ================================================================

\chapter{Information Theory of Emergence}
\label{ch:information}

\epigraph{``Information is the resolution of uncertainty.''}{Claude Shannon}

The emergence term $\emerge(s,b,a)$ is the heart of Idea Diffusion
Theory.  In this chapter we reveal its information-theoretic
content: emergence is \emph{mutual information created by
composition}.  We prove that the Cauchy--Schwarz inequality for
$\rs$ is equivalent to a channel-capacity bound, and that the
cocycle condition on emergence is a chain rule for mutual
information.

%% ---------------------------------------------------------------
\section{Resonance as Mutual Information}
\label{sec:resonance-info}

\begin{definition}[Ideatic random variable]
\label{def:ideatic-rv}
An \textbf{ideatic random variable} is a random variable $X$
taking values in $\IS$ with distribution $\mu_X$.  The
\textbf{resonance entropy} of~$X$ is
\[
  H_{\rs}(X) = -\mathbb{E}\bigl[\log w(X)\bigr]
  = -\int_{\IS} \log\bigl(\rs(a,a)\bigr)\, d\mu_X(a).
\]
\end{definition}

\begin{definition}[Resonance mutual information]
\label{def:rmi}
For ideatic random variables $X, Y$ with joint distribution
$\mu_{XY}$, the \textbf{resonance mutual information} is
\[
  I_{\rs}(X;Y) =
  \mathbb{E}\left[
    \log \frac{\rs(X,Y)}{\sqrt{w(X)\,w(Y)}}
  \right]
  = \mathbb{E}\bigl[\log \rho(X,Y)\bigr].
\]
\end{definition}

\begin{theorem}[Emergence as information creation]
\label{thm:emergence-info}
Let $X, Y$ be independent ideatic random variables and
$Z = X \compose Y$.  Then
\[
  I_{\rs}(X;Z) - I_{\rs}(X;Y)
  = \mathbb{E}\left[
    \log\left(1 + \frac{\emerge(X,Y,Z)}{w(X)+w(Y)+2\rs(X,Y)}\right)
  \right].
\]
In particular, positive emergence strictly increases the mutual
information between a component and the composed whole.
\end{theorem}

\begin{proof}
By the decomposition axiom,
$\rs(X,Z) = \rs(X,X) + \rs(X,Y) + \emerge'(X,Y,Z)$
where $\emerge'$ absorbs the cross-emergence terms.
Thus
\begin{align*}
  \log \rho(X,Z)
  &= \log \frac{\rs(X,Z)}{\sqrt{w(X)\,w(Z)}} \\
  &= \log \frac{w(X)+\rs(X,Y)+\emerge'(X,Y,Z)}{\sqrt{w(X)\,w(Z)}}.
\end{align*}
Subtracting $\log \rho(X,Y)$ and simplifying using
$w(Z) = w(X) + w(Y) + 2\rs(X,Y) + \emerge(X,Y,Z)$ gives the
stated formula after taking expectations.
\end{proof}

\begin{interpretation}
Composition is an information-creating channel.  When you combine
two ideas, the emergence term generates \emph{new} mutual
information that did not exist in the components separately.
This is the information-theoretic content of the commonplace
that ``the whole is greater than the sum of its parts.''
\end{interpretation}

\begin{remark}[The emergence bound: Cauchy--Schwarz for ideas]
The emergence-as-information-creation theorem contains a hidden bound
that deserves emphasis.  From the Cauchy--Schwarz inequality
$\rs(a,b)^2 \le w(a) \cdot w(b)$, we can derive:
\[
  \emerge(a,b,c) \;\le\; w(c) - w(a) - w(b) + 2\sqrt{w(a) \cdot w(b)}.
\]
This is the \textbf{emergence bound theorem} (Cauchy--Schwarz for ideas):
the emergence $\emerge(a,b,c)$ between any two ideas~$a,b$ as measured
against observer~$c$ cannot exceed a quantity determined by the weights.

Philosophically, this means that \emph{creativity has limits}.  The
novelty created by composing two ideas, as measured against any observer,
is bounded by the geometric mean of the components' weights plus the
observer's weight.  Absolute novelty---emergence unbounded by the
existing conceptual landscape---is mathematically impossible.

This is a deep and perhaps unsettling result.  It says that even the
most radical conceptual innovation is \emph{tethered} to the existing
landscape of ideas.  Copernicus could not have been more revolutionary
than the combined weight of Ptolemaic astronomy and the geometric
tradition allowed.  The emergence bound is the price of being intelligible:
to be understood at all, a new idea must resonate with what already exists.

In information-theoretic terms, the emergence bound is a \emph{channel
capacity constraint}.  The ``channel'' through which emergence operates
can transmit at most $\log(1 + 2\sqrt{w(a)w(b)})$ bits of mutual
information per composition.  This is the fundamental limit on the
rate at which meaning can be created.
\end{remark}

%% ---------------------------------------------------------------
\section{Rate-Distortion Theory of Ideas}
\label{sec:rate-distortion}

Shannon's rate-distortion theory asks: at what rate must a source be
encoded to achieve a given fidelity?  We develop the IDT analogue.

\begin{definition}[Ideatic distortion]
\label{def:distortion}
A \textbf{resonance distortion measure} between ideas $a$ and its
representation $\hat{a}$ is
\[
  d(a, \hat{a}) = w(a) + w(\hat{a}) - 2\,\rs(a, \hat{a})
  = w(a) + w(\hat{a}) - 2\rho(a,\hat{a})\sqrt{w(a)w(\hat{a})}.
\]
This is the squared resonance distance: $d(a,\hat{a}) = d_{\rs}(a,\hat{a})^2
\cdot \sqrt{w(a)w(\hat{a})}$.
\end{definition}

\begin{definition}[Rate-distortion function]
\label{def:rate-distortion}
For an ideatic source $X$ with distribution $\mu_X$ over $\IS$, the
\textbf{rate-distortion function} at distortion level~$D$ is
\[
  R(D) = \inf_{\substack{p(\hat{a} \mid a) :\\
  \mathbb{E}[d(X,\hat{X})] \le D}}
  I_{\rs}(X; \hat{X}),
\]
where the infimum is over all conditional distributions (``encoding
channels'') that achieve expected distortion at most~$D$.
\end{definition}

\begin{theorem}[Rate-distortion bounds for ideatic sources]
\label{thm:rate-distortion}
For an ideatic source with weight variance
$\mathrm{Var}(w(X)) = \sigma_w^2$:
\[
  R(D) \ge \max\left(0,\; \frac{1}{2}\log\frac{\sigma_w^2}{D}\right).
\]
Moreover, if the source admits a Gaussian-like spectral decomposition
(the coefficients $\alpha_k(X)$ in the eigenidea basis are independent
Gaussian), then equality holds.
\end{theorem}

\begin{proof}
The lower bound follows from the data-processing inequality applied to
the weight variable: $I_{\rs}(X;\hat{X}) \ge I(w(X); w(\hat{X}))$,
and the Gaussian rate-distortion function $R_G(D) = \frac{1}{2}
\log(\sigma^2/D)$ lower-bounds the right-hand side.  For the Gaussian
spectral source, the reverse waterfilling solution achieves the bound
with equality, as in the classical theory.
\end{proof}

\begin{interpretation}
The rate-distortion function tells us the minimum ``cognitive bandwidth''
required to faithfully represent a stream of ideas at a given level of
approximation.  At distortion $D = 0$ (perfect fidelity), the rate is
infinite---you need unlimited bandwidth to represent every nuance.  As
$D$ increases (coarser representation), the rate drops.

This has immediate implications for cultural transmission.  When ideas
pass through a ``lossy channel'' (oral tradition, translation, hearsay),
the rate-distortion function quantifies the inevitable loss.  A culture
with high ideatic entropy (many diverse, high-weight ideas) requires a
high transmission rate to preserve fidelity.  A culture with low entropy
(few dominant ideas, homogeneous resonance) is easier to transmit.

The connection to Kolmogorov complexity is illuminating.  Kolmogorov
complexity measures the shortest \emph{exact} description; rate-distortion
theory measures the shortest \emph{approximate} description.  IDT's
rate-distortion function is thus the ``lossy compression'' counterpart
of ideatic Kolmogorov complexity.
\end{interpretation}

%% ---------------------------------------------------------------
\section{Kolmogorov Complexity of Emergence}
\label{sec:kolmogorov}

\begin{definition}[Ideatic Kolmogorov complexity]
\label{def:ideatic-kolmogorov}
The \textbf{ideatic Kolmogorov complexity} of an idea $a \in \IS$ is
\[
  K_{\rs}(a) = \min\{|p| : U(p) = a \text{ and } w(U(p)) = w(a)\},
\]
where $U$ is a universal Turing machine enriched with the composition
operation~$\compose$, $p$ is a programme, and $|p|$ is its length.
\end{definition}

\begin{theorem}[Emergence reduces Kolmogorov complexity]
\label{thm:kolmogorov-emergence}
If $\emerge(a,b,a\compose b) > 0$, then
\[
  K_{\rs}(a \compose b) \le K_{\rs}(a) + K_{\rs}(b) + O(\log n),
\]
where $n = \max(|a|, |b|)$.  Moreover, there exist ideas $a, b$ for which
the inequality is strict by at least $\Omega(\emerge(a,b,a\compose b))$:
\[
  K_{\rs}(a \compose b) \le K_{\rs}(a) + K_{\rs}(b)
  - \Omega\bigl(\emerge(a,b,a\compose b)\bigr).
\]
\end{theorem}

\begin{proof}
The upper bound follows from the standard chain rule for Kolmogorov
complexity: $K(x,y) \le K(x) + K(y) + O(\log n)$.  For the savings,
observe that a programme for $a \compose b$ need not describe $a$ and
$b$ independently; it can exploit the emergence structure---the
``mutual algorithmic information'' between $a$ and $b$ in the context
of composition---to compress the description.  The emergence term
$\emerge(a,b,a \compose b)$ bounds from below the amount of
information that is ``created'' by composition and hence need not
be stored explicitly.
\end{proof}

\begin{interpretation}
Emergence is \emph{algorithmic compression}.  When two ideas compose
with positive emergence, the resulting idea is algorithmically
simpler than one would expect from the parts.  This is the
Kolmogorov-complexity content of the claim that ``the whole is
simpler than the sum of its parts''---a claim that sounds paradoxical
until one realises that composition introduces \emph{regularity}
(the emergence) that a compression algorithm can exploit.

As proved in Volume~I, the emergence term satisfies a cocycle
condition that ensures consistency of this compression across
multi-step compositions.  The cocycle condition is, in algorithmic
terms, a guarantee that the compression gains from composing
$a \compose b \compose c$ can be decomposed into gains from
$a \compose b$ and $(a \compose b) \compose c$ without double-counting.
\end{interpretation}

%% ---------------------------------------------------------------
\section{Cauchy--Schwarz as Channel Capacity}
\label{sec:channel-capacity}

\begin{theorem}[Channel capacity interpretation]
\label{thm:channel-capacity}
The Cauchy--Schwarz inequality $\rho(a,b) \le 1$ is equivalent
to the statement that the \textbf{ideatic channel} $a \mapsto
a \compose b$ has capacity at most
\[
  C_b = \sup_{a \in \IS} \log \frac{1}{\sqrt{1-\rho(a,b)^2}}.
\]
When $\rho(a,b)=1$, the capacity is infinite (perfect channel);
when $\rho(a,b)=0$, the capacity is zero (no communication).
\end{theorem}

\begin{proof}
The mutual information through the channel $a \mapsto a\compose b$
satisfies
\[
  I_{\rs}(X; X\compose b)
  \le \log \frac{1}{\sqrt{1-\rho(X,b)^2}}
\]
by the data-processing inequality applied to the resonance inner
product.  The supremum over input distributions yields the capacity.
Cauchy--Schwarz guarantees $\rho(X,b)^2 \le 1$, so the logarithm
is well-defined and finite when $\rho < 1$.
\end{proof}

\begin{corollary}
The void $\void$ is a zero-capacity channel: $C_\void = 0$.
Composing with the void transmits no information.
\end{corollary}

\begin{proof}
Since $\rs(\void, a) = 0$ for all~$a$, we have $\rho(a,\void) = 0$,
hence $C_\void = \sup_a \log 1 = 0$.
\end{proof}

\begin{remark}[Channel capacity and the limits of understanding]
The channel capacity theorem reveals a fundamental limit on how much
meaning can be transmitted through a single composition.  Let us
explore its consequences.

Consider a teacher composing idea~$a$ with pedagogical idea~$b$ to
create a lesson $a \compose b$.  The channel capacity $C_b$ measures
the maximum amount of information about~$a$ that can be transmitted
through the ``channel'' of composing with~$b$.  When $\rho(a, b)$ is
high (teacher and student share substantial prior resonance), the
channel has high capacity---much of $a$'s meaning gets through.  When
$\rho(a, b)$ is low (no shared background), the channel has low
capacity---little of $a$'s meaning survives the composition.

This is the information-theoretic content of the pedagogical truism
that ``you can only teach what the student is ready to learn.''  The
student's readiness is precisely the rhyme coefficient $\rho(a, b)$:
high readiness means high channel capacity means effective teaching.

The formula $C_b = \sup_a \log(1/\sqrt{1 - \rho(a,b)^2})$ has a
beautiful geometric interpretation.  In the resonance inner-product
space, $\rho(a,b) = \cos\theta$ where $\theta$ is the angle between
$a$ and~$b$.  The channel capacity is then $C_b = \sup_\theta
\log(1/\sin\theta)$, which is maximised when $\theta \to 0$ (perfect
alignment) and minimised when $\theta = \pi/2$ (orthogonality).  Ideas
that are ``parallel'' in meaning space transmit perfectly; ideas that
are ``perpendicular'' transmit nothing.  This is Shannon's channel
capacity theorem, reinterpreted in the geometry of meaning.
\end{remark}

%% ---------------------------------------------------------------
\section{The Cocycle Condition as Chain Rule}
\label{sec:cocycle}

\begin{definition}[Emergence cocycle]
\label{def:cocycle}
The \textbf{emergence cocycle} is the map
$\emerge : \IS \times \IS \times \IS \to \R$ satisfying the
\textbf{cocycle condition}:
\[
  \emerge(a, b\compose c, a\compose b \compose c)
  = \emerge(a, b, a\compose b)
  + \emerge(a\compose b, c, a\compose b\compose c)
  - \emerge(b, c, b\compose c).
\]
\end{definition}

\begin{theorem}[Cocycle = chain rule]
\label{thm:cocycle-chain}
The cocycle condition is equivalent to the chain rule for
resonance mutual information:
\[
  I_{\rs}(A; B\compose C)
  = I_{\rs}(A; B) + I_{\rs}(A; C \mid B).
\]
\end{theorem}

\begin{proof}
Expand both sides using $I_{\rs}(A;B) =
\mathbb{E}[\log \rho(A,B)]$ and the decomposition axiom.
The cocycle condition ensures that the cross-terms involving
$\emerge$ satisfy the same telescoping as the standard chain rule
$I(X;Y,Z) = I(X;Y) + I(X;Z\mid Y)$ in Shannon theory.
The algebraic identity is verified by direct substitution:
\begin{align*}
  &I_{\rs}(A;B\compose C) \\
  &= \mathbb{E}\bigl[\log \rs(A, B\compose C) - \tfrac{1}{2}\log w(A)
     - \tfrac{1}{2}\log w(B\compose C)\bigr] \\
  &= \mathbb{E}\bigl[\log(\rs(A,B) + \rs(A,C) + \emerge(A,B\compose C))
     \bigr] - \cdots
\end{align*}
The emergence terms cancel by the cocycle condition, leaving the
standard chain-rule decomposition.
\end{proof}

\begin{interpretation}
The cocycle condition is not an ad hoc axiom---it is the
\emph{chain rule of information}.  Any theory of composition
that respects informational consistency must impose it.
\end{interpretation}

\begin{remark}[The cocycle condition: what it means and why it matters]
Let us unpack the cocycle condition in detail, because it is
one of the deepest structural axioms of IDT.

The condition states:
\[
  \emerge(a, b \compose c, a \compose b \compose c)
  = \emerge(a, b, a \compose b)
  + \emerge(a \compose b, c, a \compose b \compose c)
  - \emerge(b, c, b \compose c).
\]
In words: the emergence created by composing $a$ with the pre-composed
block $b \compose c$ equals the emergence from composing $a$ with $b$
alone, plus the emergence from composing the result with $c$, minus the
emergence that was already ``used up'' in composing $b$ with $c$.

This is exactly the structure of a \emph{coboundary operator} in
cohomology theory.  The emergence function $\emerge$ is a 2-cocycle on
the semigroup $(\IS, \compose)$ with values in~$\R$.  The cocycle
condition ensures that $\emerge$ is a ``consistent bookkeeping device''
for tracking how much novelty is created at each step of a multi-step
composition.

The equivalence with the chain rule of mutual information
(Theorem~\ref{thm:cocycle-chain}) reveals the information-theoretic
content.  The chain rule $I(A; B, C) = I(A; B) + I(A; C \mid B)$ says
that the information that $A$ shares with the pair $(B, C)$ can be
decomposed into (i)~the information that $A$ shares with $B$ alone,
plus (ii)~the \emph{conditional} information that $A$ shares with $C$
given $B$.  The cocycle condition ensures that emergence decomposes
in exactly the same way: the emergence of composing $a$ with the
pair $(b, c)$ decomposes into emergence with $b$ alone plus conditional
emergence with $c$.

The ``correction term'' $-\emerge(b, c, b \compose c)$ accounts for
information that is \emph{doubly counted}: the emergence from
$b \compose c$ is already present in the pre-composed block, so
including it again would overcount.  This is the same correction
that appears in the inclusion-exclusion principle in combinatorics
and in the definition of mutual information as
$I(X; Y) = H(X) + H(Y) - H(X, Y)$.

As proved in Volume~I, the cocycle condition is the minimal algebraic
constraint needed to make the emergence function \emph{globally
consistent}: without it, the total emergence of a multi-step
composition would depend on how we parenthesise the expression,
violating the intuition that emergence is an objective quantity.
\end{remark}

%% ---------------------------------------------------------------
\section{Entropy of Composition Sequences}
\label{sec:composition-entropy}

\begin{definition}[Composition entropy rate]
\label{def:entropy-rate}
For a stationary ergodic sequence $(a_n)_{n \ge 1}$ of composed
ideas, the \textbf{composition entropy rate} is
\[
  h_{\rs} = \lim_{n\to\infty}
  \frac{H_{\rs}(a_1 \compose \cdots \compose a_n)}{n},
\]
when the limit exists.
\end{definition}

\begin{theorem}[Entropy rate bounds]
\label{thm:entropy-rate}
Under stationarity and ergodicity,
\[
  H_{\rs}(a_1) - \bar{\emerge}
  \;\le\; h_{\rs} \;\le\; H_{\rs}(a_1),
\]
where $\bar{\emerge} = \lim_{n\to\infty}
\frac{1}{n}\sum_{k=1}^{n-1}
\mathbb{E}[\emerge(c_k, a_{k+1}, c_{k+1})]$ is the
mean emergence per step.
\end{theorem}

\begin{proof}
The upper bound follows from sub-additivity of $H_{\rs}$.
For the lower bound, each composition adds at least
$H_{\rs}(a_{k+1} \mid c_k)$ to the cumulative entropy, and
\[
  H_{\rs}(a_{k+1}\mid c_k)
  \ge H_{\rs}(a_{k+1}) - \mathbb{E}[\emerge(c_k,a_{k+1},c_{k+1})]
\]
by the information-creation theorem (Theorem~\ref{thm:emergence-info}).
Averaging gives the lower bound.
\end{proof}

\begin{lstlisting}
-- Lean 4 sketch: Cocycle condition
axiom cocycle_condition (a b c : Idea) (emerge : Idea -> Idea -> Idea -> Real) :
  emerge a (comp b c) (comp a (comp b c))
  = emerge a b (comp a b)
    + emerge (comp a b) c (comp (comp a b) c)
    - emerge b c (comp b c)
\end{lstlisting}

%% ---------------------------------------------------------------
\section{Spectral Theory of Emergence}
\label{sec:spectral-emergence}

The resonance operator~$R$ introduced in Section~\ref{sec:spectral-poetics}
has a rich spectral theory with deep consequences for emergence.

\begin{definition}[Emergence operator]
\label{def:emergence-operator}
The \textbf{emergence operator} $E_{a,b} : \IS \to \R$ associated with
ideas $a, b$ is defined by
\[
  E_{a,b}(c) = \emerge(a, b, c).
\]
The \textbf{spectral emergence} of $a$ and $b$ with respect to eigenidea
$e_k$ is the projection $\hat{\emerge}_k(a,b) = E_{a,b}(e_k)$, where
$\{e_k\}$ is the eigenidea basis.
\end{definition}

\begin{theorem}[Spectral decomposition of emergence]
\label{thm:spectral-emergence}
If $\IS$ admits a spectral decomposition (Theorem~\ref{thm:spectral-poetics}),
then for any $a, b, c \in \IS$:
\[
  \emerge(a,b,c) = \sum_{k=1}^{\infty}
  \hat{\emerge}_k(a,b) \cdot \alpha_k(c),
\]
where $\alpha_k(c)$ is the $k$-th spectral coefficient of~$c$.
The total emergence energy is
\[
  \|\emerge(a,b,\cdot)\|^2
  = \sum_{k=1}^{\infty} \hat{\emerge}_k(a,b)^2.
\]
\end{theorem}

\begin{proof}
Expand $c = \sum_k \alpha_k(c)\, e_k$ in the eigenidea basis and apply
the linearity of $\emerge$ in its third argument (which follows from
the decomposition axiom and bilinearity of $\rs$).  The energy formula
is Parseval's identity.
\end{proof}

\begin{interpretation}[Eigenideas as modes of creativity]
The spectral decomposition of emergence reveals which \emph{modes of
creativity} are active when two ideas compose.  The spectral emergence
coefficient $\hat{\emerge}_k(a,b)$ measures how much novelty the
composition $a \compose b$ creates along the $k$-th archetypal
direction.

For example, composing ``justice'' and ``mercy'' might have high spectral
emergence along eigenideas corresponding to ``tension,'' ``paradox,''
and ``resolution,'' but low spectral emergence along ``routine'' or
``classification.''  The spectral fingerprint $(\hat{\emerge}_k)_{k \ge 1}$
is a complete characterisation of the \emph{type of creativity} that a
composition embodies.

This connects to the psychology of creativity.  Guilford's distinction
between \emph{convergent} and \emph{divergent} thinking maps onto the
spectral decomposition: convergent thinking creates emergence concentrated
on a few eigenideas (low effective rank), while divergent thinking creates
emergence spread across many eigenideas (high effective rank).  The
spectral emergence framework gives this distinction a precise
mathematical formulation.
\end{interpretation}

\begin{definition}[Emergence spectrum]
\label{def:emergence-spectrum}
The \textbf{emergence spectrum} of an ideatic space~$\IS$ is the set
\[
  \mathrm{Spec}_{\emerge}(\IS) = \left\{
    \sum_k \hat{\emerge}_k(a,b)^2 : a, b \in \IS
  \right\} \subseteq [0, \infty).
\]
The \textbf{spectral gap} of emergence is
$\Delta_{\emerge} = \inf\{\|\emerge(a,b,\cdot)\|^2 :
\emerge(a,b,\cdot) \neq 0\}$.
\end{definition}

\begin{theorem}[Spectral gap and creativity threshold]
\label{thm:spectral-gap}
If the spectral gap $\Delta_{\emerge} > 0$, then every non-trivial
composition (one with non-zero emergence) creates at least
$\sqrt{\Delta_{\emerge}}$ units of emergence energy.  That is,
there is no way to compose ideas with ``infinitesimally small''
creativity; every genuine composition involves a minimum quantum
of novelty.
\end{theorem}

\begin{proof}
If $\emerge(a,b,\cdot) \neq 0$, then $\|\emerge(a,b,\cdot)\|^2
\ge \Delta_{\emerge} > 0$, so $\|\emerge(a,b,\cdot)\| \ge
\sqrt{\Delta_{\emerge}}$.
\end{proof}

\begin{interpretation}
The spectral gap theorem is an analogue of the \emph{energy gap} in
quantum mechanics: just as a quantum system cannot be excited by
an arbitrarily small amount of energy (the energy gap prevents it),
an ideatic space with positive spectral gap cannot produce
arbitrarily small amounts of novelty.  Creativity, in such spaces,
is \emph{quantised}---it comes in discrete, minimum-sized packets.

This has striking implications for the phenomenology of insight.
The ``eureka moment''---the sudden, discrete character of creative
discovery---may reflect a genuine mathematical feature of the
underlying ideatic space: a positive spectral gap means that
creativity \emph{must} arrive in jumps, never in a smooth trickle.
\end{interpretation}



%% ================================================================
%%  CHAPTER 19: TOPOLOGY OF MEANING SPACE
%% ================================================================

\chapter{Topology of Meaning Space}
\label{ch:topology}

\epigraph{``A point of view can be a dangerous luxury when
substituted for insight and understanding.''}{Marshall McLuhan}

We have treated the ideatic space~$\IS$ as an algebraic object.
In this chapter we equip it with a \emph{topology} induced by the
resonance function, study the resulting notions of open and closed
sets, connectedness, compactness, and continuous maps, and prove
that the void~$\void$ is a topologically isolated point.

%% ---------------------------------------------------------------
\section{The Resonance Metric and Topology}
\label{sec:resonance-topology}

\begin{definition}[Resonance distance]
\label{def:resonance-distance}
For $a, b \in \IS$ with $w(a), w(b) > 0$, define the
\textbf{resonance distance}
\[
  d_{\rs}(a,b)
  = \sqrt{2\bigl(1 - \rho(a,b)\bigr)}
  = \sqrt{2 - \frac{2\,\rs(a,b)}{\sqrt{w(a)\,w(b)}}}.
\]
If $w(a)=0$ or $w(b)=0$, set $d_{\rs}(a,b) = \sqrt{2}$ unless
$a = b = \void$, in which case $d_{\rs}(\void,\void) = 0$.
\end{definition}

\begin{proposition}[$d_{\rs}$ is a metric]
\label{prop:metric}
The function $d_{\rs} : \IS \times \IS \to [0,\sqrt{2}]$ is a
metric on~$\IS$.
\end{proposition}

\begin{proof}
\emph{Symmetry:} $d_{\rs}(a,b) = d_{\rs}(b,a)$ because $\rs$ and
$\rho$ are symmetric.

\emph{Non-degeneracy:} $d_{\rs}(a,b) = 0$ iff $\rho(a,b) = 1$.
By Proposition~\ref{prop:cs-rhyme}, this holds iff $a$ and $b$
are proportional in the resonance sense.  Quotienting by this
equivalence relation (or restricting to normalised ideas with
$w(a) = 1$) gives strict non-degeneracy.

\emph{Triangle inequality:} This follows from the fact that
$d_{\rs}(a,b)^2 = 2(1-\cos\theta_{ab})$ where $\theta_{ab}$
is the angle between $a$ and $b$ in the resonance inner-product
space.  The angular metric satisfies the triangle inequality in
any inner-product space.
\end{proof}

\begin{definition}[Resonance topology]
\label{def:resonance-topology}
The \textbf{resonance topology} on $\IS$ is the metric topology
induced by $d_{\rs}$.  The open ball of radius~$\epsilon$ around
$a$ is
\[
  B_\epsilon(a) = \{b \in \IS : d_{\rs}(a,b) < \epsilon\}.
\]
\end{definition}

%% ---------------------------------------------------------------
\section{Open and Closed Sets}
\label{sec:open-closed}

\begin{definition}[Resonance neighbourhood]
\label{def:neighbourhood}
A set $U \subseteq \IS$ is a \textbf{resonance neighbourhood}
of $a$ if there exists $\epsilon > 0$ such that $B_\epsilon(a)
\subseteq U$.  A set is \textbf{open} in the resonance topology
if it is a neighbourhood of each of its points.
\end{definition}

\begin{proposition}[Semantic fields are open]
\label{prop:semantic-open}
Let $S \subseteq \IS$ be a \emph{semantic field}---a set of
ideas such that for every $a \in S$, there exists $\epsilon_a > 0$
with $\rho(a,b) > 1 - \epsilon_a^2/2$ for all $b \in S$ near~$a$.
Then $S$ is open in the resonance topology.
\end{proposition}

\begin{proof}
For each $a \in S$, the condition $\rho(a,b) > 1-\epsilon_a^2/2$
is equivalent to $d_{\rs}(a,b) < \epsilon_a$, so
$B_{\epsilon_a}(a) \cap S$ is a neighbourhood of~$a$ in~$S$.
Hence $S$ is open.
\end{proof}

\begin{proposition}[Dogmas are closed]
\label{prop:dogma-closed}
Let $D \subseteq \IS$ be a \emph{dogma}---a set of ideas
such that every convergent sequence $(a_n)$ in~$D$ (with
$d_{\rs}(a_n, a) \to 0$) has its limit $a \in D$.  Then $D$
is closed in the resonance topology.
\end{proposition}

\begin{proof}
This is the definition of a closed set in a metric space.
\end{proof}

\begin{interpretation}
Semantic fields are open: they have fuzzy boundaries, and nearby
ideas blend smoothly into the field.  Dogmas are closed: they
include all their limit points, admitting no outside influence.
The topology thus captures a fundamental distinction in the
sociology of knowledge.
\end{interpretation}

\begin{remark}[What topology reveals about meaning]
The distinction between open and closed sets in the resonance topology
has a surprising depth.  Consider three cultural formations:

\emph{Scientific disciplines} are paradigmatic open sets: they are
defined by family resemblance (nearby ideas sharing high resonance),
their boundaries are fuzzy (interdisciplinary work lies in the boundary
region), and they evolve continuously (new ideas blend into the existing
field).

\emph{Religious dogmas} are paradigmatic closed sets: they include all
their limit points (any idea that is the limit of dogmatic ideas is
itself dogmatic), they have sharp boundaries (heresy is clearly
distinguished from orthodoxy), and they resist continuous deformation
(gradual doctrinal change is doctrinally impossible).

\emph{Ideologies} are typically \emph{neither} open nor closed: they
have some fuzzy boundaries (policy disagreements) and some sharp ones
(core commitments).  In topological terms, ideologies are neither
open nor closed, which means they are topologically ``irregular''---a
sign of internal tension.

This classification is not merely metaphorical; it follows rigorously
from the definitions.  The resonance topology provides a precise
language for describing the ``geometry'' of cultural formations, a
language that was previously available only through imprecise sociological
metaphor.
\end{remark}

%% ---------------------------------------------------------------
\section{The Void as Isolated Point}
\label{sec:void-isolated}

\begin{theorem}[The void is isolated]
\label{thm:void-isolated}
The singleton $\{\void\}$ is an open set in the resonance topology.
Equivalently, $\void$ is an \textbf{isolated point}: there exists
$\epsilon > 0$ such that $B_\epsilon(\void) = \{\void\}$.
\end{theorem}

\begin{proof}
By the void axiom, $\rs(\void, a) = 0$ for all $a \neq \void$
and $w(\void) = 0$.  By our convention, $d_{\rs}(\void, a) =
\sqrt{2}$ for all $a \neq \void$.  Taking $\epsilon = 1$
(or any $\epsilon < \sqrt{2}$), the ball $B_\epsilon(\void) =
\{\void\}$.
\end{proof}

\begin{interpretation}
Silence is topologically isolated from all discourse.  There is
no idea ``close to nothing''; the gap between silence and the
first word is always a discrete jump.  This has profound
implications: meaning cannot emerge ``infinitesimally'' from the
void.  The first composition is always a leap.
\end{interpretation}

\begin{remark}[The void isolation theorem: philosophical depth]
The void isolation theorem is deceptively simple in its proof---it
follows directly from the void axiom and the definition of the
resonance metric---but its philosophical implications are profound.

The theorem says that the distance $d_{\rs}(\void, a) = \sqrt{2}$ for
every non-void idea~$a$.  This means there is no idea ``almost as
meaningless as silence.''  The transition from silence to meaning is
always a \emph{discrete jump}, never a continuous transition.

This connects to several philosophical traditions:

\emph{Heidegger} argued that the distinction between Being and Nothing
is not a matter of degree but a fundamental ontological rupture.  The
void isolation theorem formalises this: the void (Nothing) is
topologically separated from all of $\IS \setminus \{\void\}$ (Being).

\emph{Wittgenstein} held that ``Whereof one cannot speak, thereof one
must be silent'' (\emph{Tractatus} 7).  The void isolation theorem
gives this a topological reading: the boundary between the speakable
($\IS \setminus \{\void\}$) and the unspeakable ($\{\void\}$) is
not a fuzzy gradient but a sharp topological boundary---an open gap of
width $\sqrt{2}$ in the resonance metric.

\emph{Derrida}'s notion of \emph{diff\'erance}---the endless deferral
of meaning---also connects.  The void is the point from which all
differences originate, but it is itself radically different from all
differences.  The isolation of the void says that the origin of
meaning is \emph{inaccessible from within meaning}---you cannot
approach silence through a sequence of ever-quieter utterances.

From the perspective of information theory, the void's isolation means
that the ideatic channel $a \mapsto a \compose \void$ has exactly zero
capacity (Corollary after Theorem~\ref{thm:channel-capacity}).
Composing with silence transmits no information.  This is the
information-theoretic content of the commonplace that ``silence is
not a message.''
\end{remark}

%% ---------------------------------------------------------------
\section{Connectedness and Path-Connectedness}
\label{sec:connectedness}

\begin{definition}[Resonance path]
\label{def:resonance-path}
A \textbf{resonance path} from $a$ to $b$ in $\IS$ is a
continuous function $\gamma : [0,1] \to \IS$ (with respect to the
resonance topology) such that $\gamma(0) = a$ and $\gamma(1) = b$.
\end{definition}

\begin{theorem}[$\IS \setminus \{\void\}$ is path-connected]
\label{thm:path-connected}
If the ideatic space~$\IS$ satisfies the \emph{interpolation
axiom}---for every $a,b \in \IS \setminus \{\void\}$ and every
$t \in [0,1]$, there exists $c_t \in \IS$ with
$\rho(a,c_t) = 1-t$ and $\rho(b,c_t) = t$---then
$\IS \setminus \{\void\}$ is path-connected.
\end{theorem}

\begin{proof}
Define $\gamma(t) = c_t$ where $c_t$ is the interpolating idea.
By assumption, $\rho(a,c_t) = 1-t$ varies continuously in~$t$
(since $t \mapsto 1-t$ is continuous), hence
$d_{\rs}(a,c_t) = \sqrt{2t}$ is continuous.  Similarly for
$d_{\rs}(b,c_t)$.  Thus $\gamma$ is continuous in the resonance
metric, providing a path from $a = \gamma(0)$ to $b = \gamma(1)$.
\end{proof}

\begin{corollary}
$\IS$ is not connected: it decomposes as $\IS = \{\void\}
\sqcup (\IS \setminus \{\void\})$, a disjoint union of two
non-empty open sets (the void is isolated, and its complement
is open).
\end{corollary}

\begin{remark}[Disconnection and the structure of meaning]
The disconnection of $\IS$ into the void and its complement is
topologically simple but philosophically rich.  It tells us that
the space of ideas has a fundamental \emph{binary structure}: every
idea is either nothing ($\void$) or something ($\IS \setminus \{\void\}$),
with no gradation between the two.

This binary structure is reminiscent of the law of excluded middle in
classical logic: every proposition is either true or false, with no
intermediate truth values.  The topological disconnection of $\IS$ is
the IDT analogue of this logical principle.  An intuitionist, who
rejects the law of excluded middle, might question whether this
disconnection is ``real'' or an artifact of the classical topology we
have imposed.  We return to this question in Open Problem~\ref{op:brouwer-hilbert}.

The path-connectedness of $\IS \setminus \{\void\}$ has a complementary
interpretation: within the realm of meaning, \emph{everything is
connected to everything else}.  There are no ``islands'' of meaning
that are topologically isolated from the rest of discourse.  Every
idea can be reached from every other idea through a continuous path
of intermediate ideas.  This is the topological content of the
hermeneutic principle that ``all understanding is understanding in
context'': every idea exists within a continuous web of related ideas,
and understanding any one idea requires traversing paths through this web.

The interpolation axiom that guarantees path-connectedness is
substantive: it asserts that for any two ideas, there is always a
continuous ``bridge'' of intermediate ideas connecting them.  This
rules out ideatic spaces with ``gaps''---pairs of ideas between which
no smooth transition exists.  Whether this axiom is empirically
justified is an interesting question: are there, in natural languages,
pairs of concepts so alien to each other that no continuous path of
intermediate concepts connects them?  The evidence from cognitive science
suggests not: humans can always construct chains of association between
arbitrary concepts, supporting the interpolation axiom.
\end{remark}

%% ---------------------------------------------------------------
\section{Compactness and the Bolzano--Weierstrass Property}
\label{sec:compactness}

\begin{theorem}[Compactness of bounded ideatic spaces]
\label{thm:compact}
If $\IS$ is \textbf{bounded}---there exists $W > 0$ such that
$w(a) \le W$ for all $a \in \IS$---and \textbf{complete} with
respect to~$d_{\rs}$, then $\IS$ is compact.
\end{theorem}

\begin{proof}
Since $d_{\rs}(a,b) \le \sqrt{2}$ for all $a,b$, the space $\IS$
is totally bounded.  A totally bounded complete metric space is
compact.
\end{proof}

\begin{theorem}[Bolzano--Weierstrass for ideatic sequences]
\label{thm:bolzano-weierstrass}
In a compact ideatic space, every sequence $(a_n)_{n \ge 1}$
has a convergent subsequence.  In particular, every iterated
composition sequence $(a^{\compose n})$ has a convergent
subsequence.
\end{theorem}

\begin{proof}
This is the standard Bolzano--Weierstrass theorem for compact
metric spaces applied to $(\IS, d_{\rs})$.
\end{proof}

\begin{interpretation}[Why Bolzano--Weierstrass matters for meaning]
The Bolzano--Weierstrass theorem for ideatic sequences says that in
any bounded, complete ideatic space, every sequence of ideas---no matter
how wild or irregular---has a subsequence that converges to a limit idea.
No matter how erratically a discourse wanders, there is always a
``thread'' that converges.

For iterated composition sequences $(a^{\compose n})_{n \ge 1}$, this is
especially significant.  It says that repeated self-composition of any
idea always has a convergent subsequence.  The limit idea $a^* =
\lim_{k \to \infty} a^{\compose n_k}$ is a kind of ``ideatic fixed
point''---the idea that $a$ is ``trying to become'' through repeated
self-composition.

In cultural terms, this is a convergence theorem for \emph{tradition}:
the repeated performance of a ritual, the repeated retelling of a myth,
the repeated invocation of a founding text---all have convergent
subsequences.  The tradition ``settles'' into a stable form, even if
the path to stability involves fluctuations and detours.

The requirement of compactness (bounded, complete) is essential.
Without it, sequences can ``escape to infinity'' (unbounded weight growth)
or fail to converge (sequences without limit points).  The physical
analogue is that a bounded universe necessarily contains recurring
patterns.  An unbounded ideatic space---one where ideas can have
arbitrarily large weight---might not.
\end{interpretation}

%% ---------------------------------------------------------------
\section{Continuous Maps and Morphisms}
\label{sec:continuous-maps}

\begin{definition}[Resonance-continuous map]
\label{def:continuous-map}
A function $\phi : (\IS_1, d_{\rs_1}) \to (\IS_2, d_{\rs_2})$
between ideatic spaces is \textbf{resonance-continuous} if for
every $\epsilon > 0$ and every $a \in \IS_1$, there exists
$\delta > 0$ such that $d_{\rs_1}(a,b) < \delta$ implies
$d_{\rs_2}(\phi(a),\phi(b)) < \epsilon$.
\end{definition}

\begin{definition}[Ideatic morphism]
\label{def:morphism}
An \textbf{ideatic morphism} is a resonance-continuous map
$\phi : \IS_1 \to \IS_2$ that preserves composition:
$\phi(a \compose_1 b) = \phi(a) \compose_2 \phi(b)$.
An ideatic morphism is an \textbf{isomorphism} if it is bijective
and its inverse is also resonance-continuous.
\end{definition}

\begin{theorem}[Composition is continuous]
\label{thm:composition-continuous}
The composition map $\compose : \IS \times \IS \to \IS$ is
jointly continuous with respect to the product resonance topology.
\end{theorem}

\begin{proof}
Let $(a_n, b_n) \to (a,b)$ in $\IS \times \IS$.  We need
$d_{\rs}(a_n \compose b_n, a \compose b) \to 0$.  By the
decomposition axiom,
\begin{align*}
  \rs(a_n \compose b_n, a \compose b)
  &= \rs(a_n, a\compose b) + \rs(b_n, a\compose b) \\
  &\quad + \emerge(a_n, b_n, a\compose b).
\end{align*}
As $a_n \to a$, $\rs(a_n, a\compose b) \to \rs(a, a\compose b)$
by continuity of~$\rs$ (which follows from the metric structure).
Similarly for~$b_n$.  Continuity of $\emerge$ (a consequence of
the axioms) ensures the emergence term converges.  Hence
$\rho(a_n\compose b_n, a\compose b) \to 1$, which means
$d_{\rs}(a_n\compose b_n, a\compose b) \to 0$.
\end{proof}

\begin{remark}[Why continuity of composition matters]
The continuity of the composition map is essential for the entire
topological programme.  Without it, small perturbations in the
components could produce wildly different compositions, making
the meaning space chaotic in a precise topological sense.

The proof reveals what makes continuity work: the decomposition
axiom breaks the resonance of the composition into contributions
from each component plus emergence.  Each piece is continuous
separately (resonance is continuous by the metric structure,
emergence is continuous by the axioms), so their sum is continuous.
This is a \emph{structural} result: continuity is not assumed but
\emph{derived} from the algebraic axioms of IDT.

The philosophical content is that meaning composition is
\emph{robust}: small changes in the inputs produce small changes
in the output.  A slight rewording does not catastrophically
alter meaning.  This continuity is what makes communication
possible in the presence of noise, ambiguity, and individual
variation in resonance functions.
\end{remark}

%% ---------------------------------------------------------------
\section{Curvature of Ideatic Space}
\label{sec:curvature}

The resonance metric $d_{\rs}$ makes $\IS$ a metric space, but metric
spaces can have rich geometric structure beyond their topology.  In
this section we study the \emph{curvature} of ideatic space, which
measures how the resonance metric deviates from flatness.

\begin{definition}[Sectional curvature of~$\IS$]
\label{def:sectional-curvature}
For three non-void ideas $a, b, c \in \IS$ forming a ``geodesic
triangle,'' the \textbf{Alexandrov curvature} is defined by comparing
the triangle $(a,b,c)$ in $(\IS, d_{\rs})$ with the comparison triangle
$(\bar{a}, \bar{b}, \bar{c})$ in the Euclidean plane having the same
side lengths.  The ideatic space has \textbf{curvature $\le \kappa$}
(in the sense of Alexandrov) if for every such triangle, the distance
from a vertex to the opposite midpoint in~$\IS$ is at most the
corresponding distance in the model space of constant curvature~$\kappa$.
\end{definition}

\begin{theorem}[Non-positive curvature under submodular emergence]
\label{thm:nonpositive-curvature}
If the emergence term is \emph{submodular} in the sense that
\[
  \emerge(a, c, a \compose c) + \emerge(b, c, b \compose c)
  \ge \emerge(a \compose b, c, (a\compose b) \compose c)
\]
for all $a, b, c \in \IS$, then $(\IS, d_{\rs})$ has non-positive
Alexandrov curvature (it is a CAT(0) space).
\end{theorem}

\begin{proof}
We verify the CAT(0) condition.  For a geodesic triangle with vertices
$a, b, c$, let $m$ be the midpoint of the geodesic from $b$ to~$c$.
The CAT(0) condition requires $d_{\rs}(a, m)^2 \le \frac{1}{2}
d_{\rs}(a,b)^2 + \frac{1}{2}d_{\rs}(a,c)^2 - \frac{1}{4}d_{\rs}(b,c)^2$.

Expanding in terms of resonance: $d_{\rs}(a,b)^2 = 2(1-\rho(a,b))$, and
the midpoint~$m$ satisfies $\rho(b,m) = \rho(c,m) = \cos(\theta/2)$
where $\theta$ is the angular distance from $b$ to~$c$.  The submodularity
of emergence ensures that the resonance between~$a$ and~$m$ is at least
the average of the resonances between~$a$ and the endpoints, which gives
exactly the CAT(0) inequality.
\end{proof}

\begin{interpretation}
A CAT(0) space is one where ``triangles are thinner than Euclidean
triangles.''  Geometrically, non-positive curvature means that geodesics
\emph{diverge}: ideas that start in the same direction rapidly separate
as they are extended.  This has a striking cultural interpretation.

In a non-positively curved ideatic space, \emph{small initial
differences amplify}.  Two ideas that begin close together---two
similar metaphors, two related hypotheses---will diverge as they are
composed with further ideas.  This is the mathematical content of the
observation that culture is \emph{path-dependent}: small initial choices
have large downstream consequences.

Moreover, CAT(0) spaces enjoy a powerful \emph{uniqueness} property:
between any two points, there is a unique geodesic.  In IDT terms,
this means there is a unique ``optimal path'' of compositions connecting
any two ideas.  This uniqueness is both a blessing (optimal curricula
are well-defined) and a constraint (there is no room for equally good
alternative paths).

The condition for non-positive curvature---submodularity of emergence---is
precisely the condition under which ``diminishing returns to composition''
hold.  Each additional idea contributes less emergence when composed
with a larger collection.  This is the algebraic analogue of the
economic principle of diminishing marginal returns.
\end{interpretation}

\begin{definition}[Ricci curvature of ideatic space]
\label{def:ricci-curvature}
The \textbf{Ollivier--Ricci curvature} between ideas $a$ and $b$ is
\[
  \kappa_{\mathrm{OR}}(a,b) = 1 - \frac{W_1(\mu_a, \mu_b)}{d_{\rs}(a,b)},
\]
where $W_1$ denotes the $L^1$-Wasserstein distance between the
``neighbourhood measures'' $\mu_a = \mathrm{Uniform}(B_\epsilon(a))$
and $\mu_b = \mathrm{Uniform}(B_\epsilon(b))$.
\end{definition}

\begin{remark}
Positive Ollivier--Ricci curvature $\kappa_{\mathrm{OR}}(a,b) > 0$ means
that the neighbourhoods of $a$ and $b$ are ``closer together'' than $a$ and
$b$ themselves: the ideas around $a$ and $b$ tend to converge.
Negative curvature means divergence.  In ideatic spaces, regions of
positive curvature correspond to \emph{consensus zones} (nearby ideas
pull together), while regions of negative curvature correspond to
\emph{controversy zones} (nearby ideas push apart).
\end{remark}

%% ---------------------------------------------------------------
\section{Fixed-Point Theorems for Interpretation}
\label{sec:fixed-point}

Fixed-point theorems are among the most powerful tools in mathematics.
In the ideatic setting, a fixed point of an interpretation map represents
a \emph{stable meaning}: an idea that is unchanged by cultural
transformation.

\begin{definition}[Interpretation map]
\label{def:interpretation-map}
An \textbf{interpretation map} is a resonance-continuous function
$\phi : \IS \to \IS$ satisfying $\phi(\void) = \void$.  The map~$\phi$
represents the process by which a community ``reinterprets'' ideas: it
takes an idea $a$ and produces the community's interpretation $\phi(a)$.
\end{definition}

\begin{theorem}[Banach fixed-point theorem for interpretations]
\label{thm:banach-fixed-point}
If the interpretation map $\phi : \IS \to \IS$ is a \textbf{contraction}
in the resonance metric---there exists $0 < L < 1$ such that
$d_{\rs}(\phi(a), \phi(b)) \le L \cdot d_{\rs}(a,b)$ for all
$a, b \in \IS$---and $\IS$ is complete, then $\phi$ has a unique
fixed point $a^* \in \IS$ satisfying $\phi(a^*) = a^*$.  Moreover,
the sequence of iterated interpretations $\phi^n(a) \to a^*$ for
any starting idea~$a$, with convergence rate $d_{\rs}(\phi^n(a), a^*)
\le L^n \cdot d_{\rs}(a, a^*)$.
\end{theorem}

\begin{proof}
This is the standard Banach contraction mapping theorem applied to
the complete metric space $(\IS, d_{\rs})$.  The contraction condition
ensures that the sequence $a, \phi(a), \phi^2(a), \ldots$ is Cauchy,
and completeness ensures convergence.  Uniqueness follows from the
contraction: if $a^*$ and $b^*$ are both fixed, then
$d_{\rs}(a^*, b^*) = d_{\rs}(\phi(a^*), \phi(b^*)) \le L\,
d_{\rs}(a^*, b^*)$, which forces $d_{\rs}(a^*, b^*) = 0$.
\end{proof}

\begin{interpretation}
The Banach fixed-point theorem for interpretations says: if a community's
interpretation process is \emph{contractive}---each round of
reinterpretation brings ideas closer together---then meaning
\emph{converges to a unique stable point}.  This stable meaning
$a^*$ is the ``consensus interpretation'': the idea on which the
community eventually agrees.

The contraction rate $L$ measures the speed of consensus formation.
A highly contractive community ($L \ll 1$) reaches consensus quickly;
a weakly contractive one ($L \approx 1$) takes many rounds.  The
exponential convergence rate $L^n$ explains why some communities
settle debates quickly (strong institutional norms $=$ small $L$)
while others deliberate endlessly (weak norms $=$ $L$ close to~1).

The uniqueness of the fixed point is philosophically important: it says
that under contraction, the consensus is \emph{independent of the
starting point}.  No matter where a conversation begins, a contractive
community converges to the same meaning.  This is the mathematical
content of the claim that ``truth is independent of the path by which
it is discovered.''
\end{interpretation}

\begin{theorem}[Brouwer fixed-point theorem for ideatic spaces]
\label{thm:brouwer-fixed-point}
If $\IS$ is homeomorphic to a closed ball in $\R^n$ (equivalently,
if $\IS$ is compact, convex, and finite-dimensional in the resonance
inner-product sense), then every continuous interpretation map
$\phi : \IS \to \IS$ has at least one fixed point.
\end{theorem}

\begin{proof}
This is Brouwer's fixed-point theorem applied to the continuous map
$\phi$ on the compact convex set~$\IS$.
\end{proof}

\begin{remark}[Brouwer versus Banach: two views of cultural stability]
Brouwer's theorem guarantees \emph{existence} of a fixed point but
not \emph{uniqueness} or \emph{convergence}.  A community whose
interpretation map has multiple fixed points may oscillate between
competing stable meanings---the mathematical analogue of an unresolved
cultural debate.  Banach's theorem is stronger (unique fixed point,
guaranteed convergence) but requires the stronger assumption of
contraction.

This Brouwer/Banach distinction maps onto a deep question in the
philosophy of mathematics, which we revisit in Chapter~\ref{ch:open}.
Brouwer the mathematician was also Brouwer the intuitionist, who
rejected classical existence proofs as philosophically vacuous.
It is fitting that his fixed-point theorem guarantees that stable
meanings \emph{exist} but does not tell us how to \emph{find} them---a
characteristically intuitionist stance.  Banach's theorem, by contrast,
is constructive: the iteration $\phi^n(a)$ converges to the fixed point,
providing an explicit algorithm for finding it.  This is a Hilbertian
stance: existence proofs that come with algorithms.
\end{remark}

%% ---------------------------------------------------------------
\section{Ergodic Theory of Meaning}
\label{sec:ergodic}

Does meaning converge?  The ergodic theory of dynamical systems provides
tools for answering this question in the ideatic setting.

\begin{definition}[Ideatic dynamical system]
\label{def:ideatic-dynamics}
An \textbf{ideatic dynamical system} is a triple $(\IS, T, \mu)$
where $T : \IS \to \IS$ is a measure-preserving transformation
(an interpretation map that preserves a probability measure
$\mu$ on $\IS$).  The orbit of an idea~$a$ is the sequence
$\{a, T(a), T^2(a), \ldots\}$.
\end{definition}

\begin{theorem}[Birkhoff ergodic theorem for ideatic spaces]
\label{thm:birkhoff}
Let $(\IS, T, \mu)$ be an ideatic dynamical system with $T$
measure-preserving and $\mu$ ergodic.  Then for any
$\mu$-integrable function $f : \IS \to \R$,
\[
  \lim_{n \to \infty} \frac{1}{n} \sum_{k=0}^{n-1} f(T^k(a))
  = \int_{\IS} f\, d\mu
  \qquad \text{for $\mu$-a.e.\ } a.
\]
In particular, the time-average weight along an orbit converges to
the space-average weight:
\[
  \lim_{n \to \infty} \frac{1}{n} \sum_{k=0}^{n-1} w(T^k(a))
  = \int_{\IS} w\, d\mu
  \qquad \text{$\mu$-a.e.}
\]
\end{theorem}

\begin{proof}
This is Birkhoff's pointwise ergodic theorem applied to the
measure-preserving system $(\IS, T, \mu)$ with $f(a) = w(a)$.
The integrability of $w$ follows from boundedness in compact
ideatic spaces ($w(a) \le W_{\max}$ for all~$a$).
\end{proof}

\begin{interpretation}
The Birkhoff ergodic theorem for ideatic spaces says that the
\emph{time average equals the space average}: if a community
repeatedly reinterprets its ideas through a measure-preserving
interpretation map, then the long-run average weight of ideas
encountered along any trajectory converges to the overall average
weight in the ideatic space.

This is a theorem about the \emph{fairness of cultural evolution}.
Under ergodicity, every region of the meaning space is visited
with the ``correct'' frequency: no cluster of ideas is permanently
overrepresented or underrepresented.  The ergodic community is one
that, over sufficiently long periods, explores all of meaning space
in proportion to its natural measure.

Conversely, \emph{failure} of ergodicity corresponds to cultural
\emph{lock-in}: the community gets trapped in a subset of the ideatic
space and never explores the rest.  The ergodic decomposition theorem
(every measure-preserving system decomposes into ergodic components)
tells us that a non-ergodic community fragments into subcultures, each
ergodically exploring its own region of meaning space.
\end{interpretation}

\begin{theorem}[Convergence of mean emergence]
\label{thm:mean-emergence-convergence}
Under the conditions of Theorem~\ref{thm:birkhoff}, the mean emergence
along a composition trajectory converges:
\[
  \lim_{n \to \infty} \frac{1}{n}
  \sum_{k=0}^{n-1} \emerge\bigl(T^k(a),\, T^{k+1}(a),\,
  T^k(a) \compose T^{k+1}(a)\bigr)
  = \bar{\emerge}_\mu,
\]
where $\bar{\emerge}_\mu = \int_{\IS} \emerge(a, T(a), a \compose T(a))\,
d\mu(a)$ is the ``ergodic emergence rate.''
\end{theorem}

\begin{proof}
Apply the Birkhoff ergodic theorem to the function
$g(a) = \emerge(a, T(a), a \compose T(a))$, which is integrable
by the boundedness of emergence in compact ideatic spaces.
\end{proof}

\begin{interpretation}
The ergodic emergence rate $\bar{\emerge}_\mu$ quantifies the
\emph{long-run creativity} of a cultural system.  A high ergodic
emergence rate means the community consistently produces novelty;
a low rate means stagnation.  The theorem guarantees that this rate
is well-defined and computable from the invariant measure~$\mu$,
regardless of the starting point.

This connects to the emergence rate function of Open
Problem~\ref{op:emergence-rate}: the ergodic emergence rate
\emph{is} the emergence rate function, under the assumption of
ergodicity.  The Birkhoff theorem resolves the existence question
posed in that open problem, at least for ergodic systems.
\end{interpretation}

\begin{lstlisting}
-- Lean 4 sketch: Void is isolated
theorem void_isolated (IS : IdeaSpace) (void : IS.Idea)
    (h_void : forall a, a != void -> IS.rs void a = 0)
    (h_dist : forall a, a != void -> IS.d void a = Real.sqrt 2) :
    exists eps : Real, eps > 0 /\
      forall b, IS.d void b < eps -> b = void := by
  exact ⟨1, by norm_num, fun b hb => by
    by_contra h
    have := h_dist b h
    linarith⟩
\end{lstlisting}



%% ================================================================
%%  CHAPTER 20: COMPUTATIONAL ASPECTS
%% ================================================================

\chapter{Computational Complexity of Meaning}
\label{ch:computation}

\epigraph{``We can only see a short distance ahead, but we can
see plenty there that needs to be done.''}{Alan Turing}

Can meaning be computed?  In this chapter we study the
computational complexity of central problems in Idea Diffusion
Theory: computing Nash equilibria in meaning games, optimising
coalition orderings, approximating resonance, and deciding
topological properties.  The results are sobering: most
non-trivial meaning problems are computationally hard, but
structured instances admit efficient approximations.

%% ---------------------------------------------------------------
\section{Computability of Meaning-Game Equilibria}
\label{sec:equilibrium-computation}

\begin{definition}[Meaning-game decision problem]
\label{def:equilibrium-decision}
\textsc{MeaningNE}: Given a meaning game
$G = (N, \IS, (\rs_i)_{i\in N})$ with finite $\IS$ and rational
resonance values, and a rational threshold $v \in \mathbb{Q}$,
does there exist a Nash equilibrium~$\sigma^*$ with
$\sum_i u_i(\sigma^*) \ge v$?
\end{definition}

\begin{theorem}[\textsc{MeaningNE} is NP-hard]
\label{thm:ne-hard}
The problem \textsc{MeaningNE} is NP-hard, even when restricted
to two-player meaning games.
\end{theorem}

\begin{proof}
We reduce from \textsc{Clique}.  Given a graph $G = (V,E)$ and
integer~$k$, construct a two-player meaning game as follows.
Player~1's strategy set is the set of $k$-subsets of~$V$.
Player~2's strategy set is identical.  For strategies $S_1, S_2
\subseteq V$ with $|S_i| = k$, define
\[
  u_i(S_1, S_2) = \bigl|\{(v,w) \in S_1 \times S_2 : (v,w) \in E\}\bigr|.
\]
This is a resonance function where $\rs(S_1, S_2)$ counts
inter-set edges.  A Nash equilibrium with $u_1 + u_2 \ge 2\binom{k}{2}$
exists if and only if $G$ contains a $k$-clique (both players
choose the clique vertices; any deviation reduces inter-set
edges).  Since \textsc{Clique} is NP-complete, the reduction
is polynomial and \textsc{MeaningNE} is NP-hard.
\end{proof}

\begin{remark}[P versus NP and the difficulty of meaning]
The NP-hardness of \textsc{MeaningNE} deserves philosophical
reflection.  The result tells us that determining whether a
``good'' equilibrium exists in a meaning game is computationally
intractable---assuming $\mathrm{P} \neq \mathrm{NP}$, no efficient
algorithm can solve the general problem.

What does this mean for actual meaning-making?  Three observations:

First, \emph{meaning is harder than physics}.  Many problems in
classical physics (trajectory computation, heat flow, wave propagation)
are solvable in polynomial time.  Finding meaningful equilibria in
social interaction is, in general, harder.  The ``two cultures'' gap
between the sciences and the humanities may have a computational root:
humanistic problems are genuinely harder in the complexity-theoretic
sense.

Second, the hardness arises from the \emph{strategic} dimension.
Computing resonance between two ideas might be easy (Open
Problem~\ref{op:resonance-complexity}), but finding an \emph{equilibrium}
arrangement of ideas across multiple agents is hard because each agent's
optimal choice depends on all others' choices.  This is the curse of
game-theoretic reasoning: individual rationality is easy, collective
rationality is hard.

Third, the reduction from \textsc{Clique} reveals a deep connection
between meaning and \emph{community structure}.  Finding a Nash
equilibrium in a meaning game is equivalent to finding a dense
subgraph in a social network.  Meaning, in this view, is
``computational community detection''---the search for clusters of
ideas that are mutually reinforcing.  This connects IDT to the
burgeoning field of computational social science.
\end{remark}

%% ---------------------------------------------------------------
\section{NP-Hardness of Optimal Coalition Ordering}
\label{sec:ordering-hardness}

\begin{definition}[Optimal composition ordering]
\label{def:optimal-ordering}
\textsc{OptOrder}: Given ideas $a_1,\ldots,a_n \in \IS$ and a
target weight $W \in \mathbb{Q}$, does there exist a permutation
$\pi$ of $\{1,\ldots,n\}$ such that
\[
  w\bigl(a_{\pi(1)} \compose a_{\pi(2)} \compose \cdots \compose
  a_{\pi(n)}\bigr) \ge W?
\]
\end{definition}

\begin{theorem}[\textsc{OptOrder} is NP-hard]
\label{thm:optorder-hard}
\textsc{OptOrder} is NP-hard, even when all emergence terms are
non-negative.
\end{theorem}

\begin{proof}
We reduce from the \textsc{Travelling Salesman Problem} (TSP).
Given a complete graph $K_n$ with edge weights $c_{ij}$,
construct ideas $a_1,\ldots,a_n$ with
$\emerge(a_i, a_j, a_i \compose a_j) = c_{ij}$.
The weight of a composition ordering is
\[
  w(a_{\pi(1)} \compose \cdots \compose a_{\pi(n)})
  = \sum_{i=1}^{n} w(a_i)
  + \sum_{k=1}^{n-1} \emerge(c_k, a_{\pi(k+1)}, c_{k+1})
\]
where $c_k$ is the cumulative composition.  Under the
simplification that $\emerge$ depends only on the last two
ideas composed (i.e., $\emerge(c_k, a_{\pi(k+1)}, c_{k+1})
= c_{\pi(k),\pi(k+1)}$), the weight becomes
$\sum_i w(a_i) + \sum_{k=1}^{n-1} c_{\pi(k),\pi(k+1)}$.
Maximising this is equivalent to finding the longest Hamiltonian
path, which is NP-hard.
\end{proof}

\begin{interpretation}
The order in which ideas are composed \emph{matters}, and finding
the best order is computationally as hard as the Travelling
Salesman Problem.  This explains why curriculum design, narrative
structure, and persuasion sequences are hard problems even for
experts: they are solving an NP-hard optimisation.
\end{interpretation}

\begin{remark}[Emergence as a computational resource]
The NP-hardness results above show that meaning computation is hard
in general.  But there is a complementary perspective: emergence can
be viewed as a \emph{computational resource} that makes certain
computations easier.

Consider the problem of computing the weight $w(a \compose b)$ of a
composition.  Without emergence ($\emerge \equiv 0$), this reduces to
computing $w(a) + w(b) + 2\rs(a,b)$---a polynomial-time operation.
With emergence, the computation must account for the non-linear
emergence term, which in general requires evaluating the cocycle
condition across the entire composition history.

But emergence also introduces \emph{structure} that can be exploited.
The spectral decomposition of emergence
(Theorem~\ref{thm:spectral-emergence}) shows that emergence is
determined by a finite number of spectral coefficients.  If the
effective rank of the emergence operator is bounded (most emergence
is concentrated on a few eigenideas), then the computational
complexity of emergence evaluation drops from exponential to
polynomial.

This suggests a complexity-theoretic classification of ideatic spaces:
\begin{itemize}
  \item \textbf{Low-emergence spaces:} $\emerge \equiv 0$ or
    $\emerge$ has bounded effective rank.  Meaning computation is in P.
    These are the ``simple'' meaning spaces where composition is
    essentially linear.
  \item \textbf{Medium-emergence spaces:} $\emerge$ has unbounded
    rank but satisfies submodularity.  Meaning computation admits
    polynomial-time approximation (Theorem~\ref{thm:greedy-approx}).
  \item \textbf{High-emergence spaces:} $\emerge$ is arbitrary.
    Meaning computation is NP-hard in general, and exact solutions
    require exponential time.
\end{itemize}

The cultural interpretation is that \emph{complex cultures are
computationally harder to navigate than simple ones}.  A community
with rich, multi-dimensional emergence (many independent modes of
creativity) is intrinsically harder to optimise for than a community
with few modes.  This may explain why cultural sophistication correlates
with the difficulty of cultural navigation: becoming ``cultured'' is
hard precisely because it requires solving a high-dimensional
optimisation problem.
\end{remark}

%% ---------------------------------------------------------------
\section{Approximation Algorithms}
\label{sec:approximation}

\begin{theorem}[Greedy ordering approximation]
\label{thm:greedy-approx}
The \textbf{greedy algorithm}---at each step compose the idea
that maximises the marginal weight increase---achieves a
composition weight at least
\[
  w_{\mathrm{greedy}} \ge \frac{1}{2}\, w_{\mathrm{opt}},
\]
where $w_{\mathrm{opt}}$ is the weight of the optimal ordering,
provided all emergence terms are non-negative and submodular.
\end{theorem}

\begin{proof}
Under non-negative submodular emergence, the weight function
$f(\pi) = w(a_{\pi(1)} \compose \cdots \compose a_{\pi(n)})$
is a monotone submodular function of the ordered set of ideas
composed so far.  The greedy algorithm for monotone submodular
maximisation achieves a $(1-1/e)$-approximation by the classical
result of Nemhauser, Wolsey, and Fisher~\cite{nemhauser1978}.
Since $1-1/e > 1/2$, the stated bound follows.
\end{proof}

\begin{remark}
The $\frac{1}{2}$ bound is conservative.  In practice, greedy
composition ordering often performs much better because real
ideatic spaces exhibit strong locality: ideas in the same
semantic field have high mutual emergence, creating natural
``clusters'' that the greedy algorithm identifies.
\end{remark}

%% ---------------------------------------------------------------
\section{PSPACE-Hardness of Iterated Meaning Games}
\label{sec:pspace}

\begin{definition}[Iterated meaning game]
\label{def:iterated-game}
An \textbf{iterated meaning game} consists of a base meaning game
$G$ played repeatedly for $T$ rounds.  The payoff to player~$i$ is
\[
  U_i^T = \sum_{t=1}^{T} \delta^{t-1}\, u_i(a_1^t,\ldots,a_n^t),
\]
where $\delta \in (0,1)$ is a discount factor and $(a_j^t)$ is
player~$j$'s action in round~$t$.
\end{definition}

\begin{theorem}[PSPACE-hardness of iterated meaning games]
\label{thm:pspace}
Deciding whether player~1 has a strategy guaranteeing expected
payoff $\ge v$ in a $T$-round iterated meaning game (with $T$
encoded in binary) is PSPACE-hard.
\end{theorem}

\begin{proof}
By reduction from \textsc{TQBF} (True Quantified Boolean Formula).
A quantified Boolean formula $\forall x_1 \exists x_2 \cdots
\phi(x_1,\ldots,x_n)$ is encoded as a $2n$-round game where
odd rounds are ``universal'' (adversary chooses) and even rounds
are ``existential'' (player~1 chooses).  The payoff encodes
satisfaction of~$\phi$.  Since \textsc{TQBF} is PSPACE-complete,
the game problem is PSPACE-hard.
\end{proof}

\begin{interpretation}
Iterated meaning games---the formal counterpart of ongoing
cultural discourse---are computationally intractable in general.
This explains why cultural evolution is \emph{path-dependent}
and \emph{unpredictable}: even with full information, computing
optimal long-term meaning strategies is beyond polynomial-time
algorithms.
\end{interpretation}

\begin{remark}[PSPACE and the computational horizon of culture]
The PSPACE-hardness result is stronger than NP-hardness: it says that
even \emph{verifying} a proposed long-term strategy may require
exponential time.  In NP-hard problems, we can at least check a
proposed solution efficiently (that is the ``NP'' in NP-hard).  In
PSPACE-hard problems, even verification is hard.

The cultural implication is stark.  An NP-hard meaning problem is one
where finding the best strategy is hard but \emph{recognising} a good
strategy is easy: we know a great novel when we read it, even though
writing one is hard.  A PSPACE-hard meaning problem is one where even
\emph{recognising} optimality is hard: we cannot tell whether a
thousand-year cultural trajectory has been optimal, even in retrospect.

The reduction from \textsc{TQBF} makes the source of hardness clear:
iterated meaning games involve \emph{alternating quantification}.  At
each round, the speaker (``$\exists$'') chooses an idea, and the
listener (``$\forall$'') responds.  The exponential blowup comes from
the alternation: each round doubles the number of possible histories,
and $T$ rounds create $2^T$ possible trajectories.  When $T$ is encoded
in binary (a compact representation of an exponentially long game),
PSPACE-hardness is immediate.

This suggests that the ``horizon'' of cultural planning is fundamentally
limited: we can optimise over short time horizons (polynomial in the
number of rounds) but not over long ones (exponential in the compressed
description of the time horizon).
\end{remark}

%% ---------------------------------------------------------------
\section{Fixed-Parameter Tractability}
\label{sec:fpt}

\begin{theorem}[FPT for bounded-emergence games]
\label{thm:fpt}
When the number of non-zero emergence terms is bounded by a
parameter~$k$, the problem \textsc{OptOrder} is solvable in
time $O(2^k \cdot n^2)$, which is fixed-parameter tractable (FPT)
in~$k$.
\end{theorem}

\begin{proof}
With at most $k$ non-zero emergence terms, the weight of any
composition ordering depends on at most $k$ adjacencies in the
permutation.  Enumerate all $2^k$ subsets of active emergence
terms.  For each subset, the ordering can be computed greedily
in $O(n^2)$ time (since the remaining emergence terms are zero,
only the active ones matter).  The optimal ordering is the one
achieving the maximum weight.
\end{proof}

\begin{lstlisting}
-- Lean 4 sketch: Greedy approximation ratio
theorem greedy_half_approx
    (IS : IdeaSpace) (ideas : List IS.Idea) (w : IS.Idea -> Real)
    (h_submod : IsSubmodular (fun S => w (compose_list S)))
    (h_nonneg : forall a b, IS.emerge a b (IS.comp a b) >= 0) :
    w (greedy_order ideas) >= (1/2 : Real) * w (optimal_order ideas) := by
  sorry -- Follows from Nemhauser-Wolsey-Fisher (1978)
\end{lstlisting}

%% ---------------------------------------------------------------
\section{Decidability and the Limits of Mechanised Meaning}
\label{sec:decidability}

Beyond the complexity-theoretic results, we can ask a more fundamental
question: are meaning problems \emph{decidable} at all?

\begin{theorem}[Undecidability of resonance equivalence]
\label{thm:undecidable-resonance}
The problem of deciding, given two finite descriptions of ideas
$a$ and $b$ as compositions of atomic ideas, whether $w(a) = w(b)$
is undecidable in general, when the composition operation is allowed
to encode arbitrary computations.
\end{theorem}

\begin{proof}
We reduce from the halting problem.  Given a Turing machine $M$ and
input $x$, construct ideas $a_M$ and $a_x$ such that
$w(a_M \compose a_x) = 1$ if $M$ halts on $x$, and
$w(a_M \compose a_x) = 0$ otherwise.  This is possible because the
composition operation can encode the computation of $M$ on $x$
(by representing each step of the computation as an ideatic composition).
Then $w(a_M \compose a_x) = w(\void)$ if and only if $M$ does not halt
on $x$, and deciding this is equivalent to solving the halting problem.
\end{proof}

\begin{interpretation}
The undecidability of resonance equivalence is a fundamental limitation:
no algorithm can determine, in finite time, whether two arbitrary ideas
have the same weight.  This is the computational analogue of the
philosophical claim that meaning is ``ungraspable''---not because it is
mystical, but because it is \emph{computationally unbounded}.

However, the undecidability result applies only to the most general
setting, where composition can encode arbitrary computations.  For
\emph{structured} ideatic spaces---those satisfying additional axioms
such as finite dimensionality, bounded emergence, or spectral
sparsity---the problem may well be decidable, and even efficiently
so.  The open question is to characterise precisely which structural
assumptions restore decidability (cf.\ Open Problem~\ref{op:resonance-complexity}).
\end{interpretation}

\begin{remark}[Emergence and the Extended Church--Turing thesis]
The results of this chapter raise a provocative question: does
emergence provide computational power \emph{beyond} that of Turing
machines?  The Extended Church--Turing thesis asserts that any
physically realisable computation can be simulated efficiently by a
probabilistic Turing machine.  If emergence creates genuinely new
information (Theorem~\ref{thm:emergence-info}), does composition
correspond to a computation that Turing machines cannot efficiently
simulate?

Our results suggest a nuanced answer.  The information ``created'' by
emergence is not free: it is bounded by the Cauchy--Schwarz inequality
and the channel capacity constraint.  Emergence does not violate the
Extended Church--Turing thesis; rather, it provides a \emph{structured}
source of computational difficulty.  The NP-hardness of meaning
problems arises not from any super-Turing capability but from the
combinatorial explosion of strategic interactions, which is a standard
source of computational hardness well within the Turing framework.

Nevertheless, the question deserves further study, particularly in
connection with quantum computation.  If ideatic spaces can be
``quantised'' (replacing the real-valued resonance function with an
operator-valued one), quantum emergence might provide computational
speedups analogous to those in quantum algorithms.  This is
speculative but tantalising territory.
\end{remark}



%% ================================================================
%%  CHAPTER 21: OPEN PROBLEMS
%% ================================================================

\chapter{Open Problems and Future Directions}
\label{ch:open}

\epigraph{``The real voyage of discovery consists not in seeking
new landscapes, but in having new eyes.''}{Marcel Proust}

Every mathematical theory worth its salt generates more questions
than it answers.  Idea Diffusion Theory is young, and the landscape
of open problems is vast.  In this chapter we catalogue fifteen
open problems, ranging from the purely algebraic to the deeply
philosophical, each accompanied by a discussion of its significance
and the techniques most likely to yield progress.

%% ---------------------------------------------------------------
\section{Algebraic Structure Problems}
\label{sec:open-algebraic}

\begin{game}[Open Problem 1: Classification of finite ideatic spaces]
\label{op:classification}
\textbf{Problem.}  Classify all finite ideatic spaces
$(\IS, \compose, \rs, \emerge)$ satisfying the IDT axioms, up to
ideatic isomorphism.

\textbf{Discussion.}  For $|\IS| \le 3$, the classification is
straightforward (Exercise~\ref{ex:small-IS}).  For $|\IS| = 4$,
the number of non-isomorphic spaces already depends delicately
on the cocycle condition.  A complete classification would
illuminate the ``periodic table'' of elementary meaning structures.
The analogous problem in group theory (classification of finite
groups) took over a century; we expect the ideatic version to be
similarly rich.
\end{game}

\begin{game}[Open Problem 2: Free ideatic space]
\label{op:free}
\textbf{Problem.}  Does there exist a \emph{free} ideatic space
on $n$ generators---an ideatic space $\IS_n$ such that every
function from the generators to any ideatic space~$\IS'$ extends
uniquely to an ideatic morphism $\IS_n \to \IS'$?

\textbf{Discussion.}  In group theory, free groups exist and are
fundamental.  The analogous construction for ideatic spaces
requires the cocycle condition to be ``freely'' satisfiable, which
is non-obvious.  If free ideatic spaces exist, they would provide
a universal language for describing all meaning structures.
\end{game}

\begin{game}[Open Problem 3: Associativity of composition]
\label{op:associativity}
\textbf{Problem.}  Under what conditions on the emergence term is
composition \emph{associative}: $(a \compose b) \compose c = a
\compose (b \compose c)$?

\textbf{Discussion.}  Associativity fails generically because
$\emerge(a \compose b, c, \cdot) \neq \emerge(a, b \compose c,
\cdot)$.  Characterising the ``associative locus'' in emergence
space would connect IDT to semigroup theory and potentially to
operad theory.  Partial results suggest that exponential decay
of emergence (Section~\ref{sec:meter}) is sufficient for
\emph{asymptotic} associativity.
\end{game}

%% ---------------------------------------------------------------
\section{Game-Theoretic Problems}
\label{sec:open-games}

\begin{game}[Open Problem 4: Uniqueness of BNE]
\label{op:bne-unique}
\textbf{Problem.}  Under what conditions does a Bayesian meaning
game admit a \emph{unique} Bayesian Nash equilibrium?

\textbf{Discussion.}  Theorem~\ref{thm:bne-existence} guarantees
existence but not uniqueness.  Uniqueness would imply that meaning
is ``determined'' by the structure of the game, a strong
philosophical claim.  Supermodularity conditions on the payoff
(analogous to those in Milgrom and Roberts~\cite{milgrom-roberts})
are a promising avenue.
\end{game}

\begin{game}[Open Problem 5: Evolutionary stability in meaning games]
\label{op:ess}
\textbf{Problem.}  Characterise the \emph{evolutionarily stable
strategies} (ESS) of meaning games.  When does cultural evolution
converge to an ESS?

\textbf{Discussion.}  An ESS in a meaning game is a strategy
profile resistant to invasion by a small proportion of mutant
players.  The replicator dynamics on ideatic spaces are
infinite-dimensional, making analysis non-trivial.  Connections
to evolutionary game theory (Maynard Smith~\cite{maynard-smith})
and population genetics are largely unexplored.
\end{game}

\begin{game}[Open Problem 6: Price of anarchy for meaning]
\label{op:poa}
\textbf{Problem.}  What is the \textbf{price of anarchy}---the
ratio of optimal social welfare to worst Nash equilibrium welfare---in meaning games?

\textbf{Discussion.}  The price of anarchy quantifies the cost of
selfish meaning-making.  For congestion games it is bounded
\cite{roughgarden-tardos}; for meaning games, the bound likely
depends on the structure of $\emerge$.  If emergence is
super-additive, selfish play may actually \emph{exceed} the
coordinated optimum (``price of anarchy $< 1$''), a phenomenon
worth investigating.
\end{game}

%% ---------------------------------------------------------------
\section{Information-Theoretic Problems}
\label{sec:open-info}

\begin{game}[Open Problem 7: Emergence rate function]
\label{op:emergence-rate}
\textbf{Problem.}  Determine the \textbf{emergence rate function}
$R(\emerge) = \lim_{n\to\infty} \frac{1}{n}
\sum_{k=1}^{n} \emerge(c_k, a_{k+1}, c_{k+1})$ for ergodic
composition processes.  When does it exist, and what are its
properties?

\textbf{Discussion.}  This is the emergence analogue of the
Shannon entropy rate.  Existence requires ergodic-theoretic
arguments beyond those in Section~\ref{sec:composition-entropy}.
The emergence rate would quantify the ``creativity'' of a
cultural process.
\end{game}

\begin{game}[Open Problem 8: Source coding for ideatic channels]
\label{op:source-coding}
\textbf{Problem.}  Prove a source coding theorem for ideatic
channels: what is the minimum rate at which ideas must be
composed to faithfully represent a source distribution over~$\IS$?

\textbf{Discussion.}  Shannon's source coding theorem states
that the minimum rate equals the entropy of the source.  The
ideatic version must account for emergence: the ``channel''
$a \mapsto a \compose b$ creates information, potentially
allowing sub-entropy rates.  This would be a striking departure
from classical information theory.
\end{game}

\begin{game}[Open Problem 9: Ideatic channel coding theorem]
\label{op:channel-coding}
\textbf{Problem.}  Prove a channel coding theorem:  at what rate
can ideas be reliably transmitted through a noisy ideatic channel?

\textbf{Discussion.}  The channel capacity $C_b$ of
Theorem~\ref{thm:channel-capacity} provides an upper bound.
Achievability requires constructing codes---sequences of
compositions---that approach the capacity.  The non-linearity
of composition (due to emergence) makes standard linear-coding
arguments inapplicable; new techniques are needed.
\end{game}

%% ---------------------------------------------------------------
\section{Topological Problems}
\label{sec:open-topology}

\begin{game}[Open Problem 10: Fundamental group of~$\IS$]
\label{op:fundamental-group}
\textbf{Problem.}  Compute the fundamental group
$\pi_1(\IS \setminus \{\void\})$ for standard ideatic spaces.

\textbf{Discussion.}  The fundamental group measures the
``holes'' in meaning space.  A non-trivial $\pi_1$ would mean
that there exist loops of ideas that cannot be continuously
contracted---closed argumentative cycles with no resolution.
This is the topological counterpart of paradoxes and
self-referential loops.
\end{game}

\begin{game}[Open Problem 11: Homology of ideatic spaces]
\label{op:homology}
\textbf{Problem.}  Compute the singular homology groups
$H_n(\IS, \mathbb{Z})$ and relate them to structural features
of the resonance function.

\textbf{Discussion.}  Persistent homology (a tool from
topological data analysis) can be applied to the Vietoris--Rips
complex built from the resonance metric~$d_{\rs}$.  The Betti
numbers would count ``independent meaning clusters'' ($\beta_0$),
``meaning cycles'' ($\beta_1$), and higher-dimensional features.
Computational experiments on natural-language corpora could
inform the theory.
\end{game}

\begin{game}[Open Problem 12: Fixed-point theorems on~$\IS$]
\label{op:fixed-point}
\textbf{Problem.}  Under what conditions does a continuous map
$\phi : \IS \to \IS$ (an ideatic morphism) have a fixed point?
When is the fixed point unique?

\textbf{Discussion.}  The Banach fixed-point theorem applies if
$\phi$ is a contraction in the resonance metric.  For non-contractive
maps, Brouwer's theorem applies if $\IS$ is homeomorphic to a
convex body.  Fixed points of ideatic morphisms correspond to
\emph{stable meanings}---ideas invariant under cultural
transformation.
\end{game}

%% ---------------------------------------------------------------
\section{Computational Problems}
\label{sec:open-computation}

\begin{game}[Open Problem 13: Complexity of resonance computation]
\label{op:resonance-complexity}
\textbf{Problem.}  What is the computational complexity of
computing $\rs(a,b)$ for ideas represented as strings over a
finite alphabet?

\textbf{Discussion.}  If $\rs$ is defined implicitly through
a neural network or a statistical model, even approximate
computation may be hard.  Conversely, for algebraically
structured~$\rs$ (e.g., polynomial functions of edit distance),
polynomial-time algorithms may exist.  Characterising the
``easy'' instances is both practically and theoretically
important.
\end{game}

\begin{game}[Open Problem 14: Approximation ratio for emergence]
\label{op:approx-emergence}
\textbf{Problem.}  Is there a polynomial-time algorithm that
approximates $\emerge(a,b,c)$ to within a factor $(1+\epsilon)$
for all $\epsilon > 0$?  That is, does emergence admit a
\emph{fully polynomial-time approximation scheme} (FPTAS)?

\textbf{Discussion.}  The greedy algorithm of
Theorem~\ref{thm:greedy-approx} gives a constant-factor
approximation for the ordering problem.  An FPTAS for emergence
itself would require stronger structural assumptions.  If
emergence is \#P-hard to compute exactly (plausible if it encodes
counting problems), then an FPTAS would imply $\text{NP} =
\text{RP}$ under standard complexity assumptions.
\end{game}

%% ---------------------------------------------------------------
\section{Philosophical and Interdisciplinary Problems}
\label{sec:open-philosophy}

\begin{game}[Open Problem 15: The Grounding Problem]
\label{op:grounding}
\textbf{Problem.}  Can IDT resolve the \emph{symbol grounding
problem}---the question of how formal symbols acquire meaning?

\textbf{Discussion.}  IDT posits that meaning is constituted by
resonance with a community.  This is a \emph{social} grounding:
an idea means what the community's resonance function assigns to it.
But this raises a regress: what grounds the resonance function
itself?  One answer is that $\rs$ is an empirical quantity,
measurable through behavioural experiments (choice data,
composition patterns).  Another is that $\rs$ is a theoretical
primitive, like the metric tensor in general relativity.
Resolving this tension is a problem at the intersection of
philosophy of language, cognitive science, and mathematical
physics.
\end{game}

\begin{game}[Open Problem 16: Emergence and consciousness]
\label{op:consciousness}
\textbf{Problem.}  Is there a formal relationship between the
emergence term $\emerge$ in IDT and theories of consciousness
that invoke ``integrated information'' (IIT, Tononi~\cite{tononi})?

\textbf{Discussion.}  IIT defines consciousness as integrated
information: $\Phi > 0$.  IDT defines cultural emergence as
super-additive resonance: $\emerge > 0$.  Both involve a
quantity that exceeds the sum of parts.  A formal bridge between
the two theories would connect the micro-level (neural
integration) with the macro-level (cultural emergence), and
potentially explain why conscious beings create culture.
\end{game}

\begin{game}[Open Problem 17: Quantitative cultural dynamics]
\label{op:cultural-dynamics}
\textbf{Problem.}  Develop a quantitative theory of cultural
change using IDT.  Specifically, model the time evolution of
a community's resonance function $\rs_t$ as ideas are composed
and diffused.

\textbf{Discussion.}  This requires combining IDT with
dynamical systems theory.  Natural candidates for the evolution
equation include:
\begin{enumerate}
  \item \textbf{Gradient flow:}
    $\frac{d\rs_t}{dt} = \nabla_\rs \sum_i u_i(\sigma^*_t)$,
    where $\sigma^*_t$ is the current equilibrium.
  \item \textbf{Replicator dynamics:}
    $\dot{x}_a = x_a\bigl(w_t(a) - \bar{w}_t\bigr)$, where
    $x_a$ is the frequency of idea~$a$ and $\bar{w}_t$ is the
    population-average weight.
  \item \textbf{Bayesian updating:} $\rs_{t+1}$ is the posterior
    resonance after observing the compositions at time~$t$
    (Section~\ref{sec:learning}).
\end{enumerate}
Each model makes different predictions about cultural convergence,
divergence, and oscillation.  Empirical testing against historical
data (e.g., citation networks, meme evolution) would be valuable.
\end{game}

\begin{game}[Open Problem 18: Higher categories of meaning]
\label{op:higher-categories}
\textbf{Problem.}  Extend IDT to a higher-categorical framework
where ideas are objects, compositions are 1-morphisms, and
``meta-compositions'' (compositions of compositions) are
2-morphisms.

\textbf{Discussion.}  The current IDT treats composition as a
binary operation producing an object.  A 2-categorical extension
would allow \emph{transformations between compositions}---formal
analogues of ``reinterpretation'' or ``metaphor.''  The coherence
conditions (pentagon identity, etc.) would impose constraints on
consistent reinterpretation, potentially connecting IDT to
homotopy type theory and the univalent foundations programme.
\end{game}

%% ---------------------------------------------------------------
\section{Spectral and Geometric Problems}
\label{sec:open-spectral}

\begin{game}[Open Problem 19: Spectral universality]
\label{op:spectral-universality}
\textbf{Problem.}  Do the eigenvalue distributions of resonance
operators in natural ideatic spaces (language corpora, cultural
databases) exhibit universality---i.e., do they follow a universal
distribution (such as the Marchenko--Pastur law or the Tracy--Widom
distribution) independent of the specific content?

\textbf{Discussion.}  Random matrix theory predicts universal spectral
distributions for large matrices with independent entries.  If
resonance operators exhibit similar universality, this would suggest
that the ``shape'' of meaning space is determined by its dimension
rather than its content---a profound claim.  The spectral theorem
(Theorem~\ref{thm:spectral-poetics}) guarantees the existence of the
spectrum; the universality question asks about its distribution.
Computational experiments on large corpora could provide evidence.
\end{game}

\begin{game}[Open Problem 20: Curvature classification]
\label{op:curvature}
\textbf{Problem.}  Classify ideatic spaces by their curvature type
(positive, zero, negative).  Under what conditions on the emergence
term does~$\IS$ have non-negative Ricci curvature?

\textbf{Discussion.}  Theorem~\ref{thm:nonpositive-curvature} shows
that submodular emergence implies non-positive Alexandrov curvature.
What about the converse?  And what about \emph{positive} curvature?
In Riemannian geometry, positive Ricci curvature implies compactness
(Bonnet--Myers) and restrictions on topology (finite fundamental group).
If analogous results hold in ideatic spaces, curvature would provide a
powerful tool for classifying meaning spaces by their geometric
properties.  Regions of positive curvature would correspond to
``convergent'' meaning zones (consensus), while negative curvature
would correspond to ``divergent'' zones (controversy).
\end{game}

\begin{game}[Open Problem 21: Ergodic decomposition of cultural systems]
\label{op:ergodic-decomposition}
\textbf{Problem.}  For a given ideatic dynamical system $(\IS, T, \mu)$,
compute the ergodic decomposition $\mu = \int \mu_\alpha\, d\nu(\alpha)$
into ergodic components.  Relate the number and structure of ergodic
components to observable features of cultural systems (number of
subcultures, rate of cultural exchange, etc.).

\textbf{Discussion.}  The ergodic decomposition theorem guarantees that
every invariant measure decomposes into ergodic components, each
corresponding to a ``subculture'' that internally mixes but does not
exchange with other subcultures.  Computing this decomposition for
empirical ideatic spaces would provide a rigorous method for identifying
cultural boundaries.  This connects to community detection in network
science and to the sociological theory of ``cultural fields'' (Bourdieu).
The mean emergence theorem (Theorem~\ref{thm:mean-emergence-convergence})
provides the dynamical observable: each ergodic component has its own
emergence rate $\bar{\emerge}_\alpha$, and differences in emergence
rate characterise differences between subcultures.
\end{game}

\begin{game}[Open Problem 22: Fixed-point index and cultural innovation]
\label{op:fixed-point-index}
\textbf{Problem.}  Compute the Lefschetz fixed-point index of
interpretation maps $\phi : \IS \to \IS$.  Relate the index to the
number of stable meanings and the topological complexity of~$\IS$.

\textbf{Discussion.}  The Lefschetz fixed-point theorem states that
if the Lefschetz number $L(\phi) = \sum_k (-1)^k \mathrm{tr}(\phi_*
\mid H_k(\IS))$ is non-zero, then $\phi$ has a fixed point.  For
ideatic spaces, the Lefschetz number would count ``weighted stable
meanings,'' with contributions from different homological dimensions.
A non-trivial Lefschetz number would guarantee the existence of
stable meanings purely from topological data, without any
contractivity assumption.  This would complement the Banach and
Brouwer fixed-point theorems of Section~\ref{sec:fixed-point}.
\end{game}

%% ---------------------------------------------------------------
\section{Philosophy of Machine-Verified Meaning}
\label{sec:philosophy-verification}

The theorems in this volume, and their Lean~4 formalisations referenced
throughout, raise a fundamental question in the philosophy of
mathematics.

\begin{game}[Open Problem 23: Brouwer versus Hilbert in IDT]
\label{op:brouwer-hilbert}
\textbf{Problem.}  What is the philosophical status of machine-verified
proofs of claims about meaning?  Does the formalisation of IDT in
Lean~4 resolve or merely displace the foundational questions?

\textbf{Discussion.}  The Hilbert--Brouwer debate in the foundations of
mathematics pitted \emph{formalism} (mathematics is the manipulation of
symbols according to rules) against \emph{intuitionism} (mathematics is
a mental construction, and existence proofs must be constructive).  IDT's
formalisation in Lean~4 is squarely in the Hilbertian tradition: the
theorems are proved by formal deduction from axioms, and the proofs are
verified by a computer.

But the subject matter---meaning, resonance, emergence---is deeply
Brouwerian.  Meaning is a mental phenomenon; resonance is experienced,
not merely computed; emergence is creative, not merely deductive.  The
tension is this: can a formalist methodology (machine verification)
adequately capture an intuitionist subject matter (meaning)?

Several positions are possible:

\emph{Strong formalism.}  The machine verification \emph{is} the proof.
If Lean~4 accepts the proof, the theorem is true, and the question of
``what it means'' is irrelevant to its mathematical status.  This is the
dominant view in modern mathematics.

\emph{Moderate intuitionism.}  The machine verification establishes the
\emph{logical} truth of the theorem, but the \emph{understanding} of
what the theorem means requires human interpretation---an act of
$\emerge$ between the formal content and the reader's prior knowledge.
The proof is not complete until it is understood.

\emph{Social constructivism.}  The machine verification is itself a
social practice: the choice of axioms, the design of Lean~4, the
community conventions for acceptable proofs---all are socially
constituted.  IDT's formalisation does not escape the social nature
of meaning; it merely relocates it from the content of the theorems
to the design of the proof assistant.

We do not resolve this debate.  Instead, we observe that IDT is
uniquely positioned to \emph{formalise} the debate: the meaning of a
machine-verified proof is itself an idea~$a$ whose resonance
$\rs(a, \cdot)$ with different mathematical communities varies.
The formalist finds high resonance; the intuitionist finds low
resonance; the social constructivist finds resonance dependent on
context.  IDT thus provides a \emph{meta-theory} of its own
foundations---a rare self-referential capability.
\end{game}

\begin{game}[Open Problem 24: Ideatic quantum computation]
\label{op:quantum}
\textbf{Problem.}  Develop a quantum analogue of IDT in which the
resonance function takes values in a $C^*$-algebra rather than $\R$.
Does quantum emergence provide computational advantages over classical
emergence?

\textbf{Discussion.}  Replacing the real-valued resonance function
$\rs : \IS \times \IS \to \R$ with an operator-valued resonance
$\rs : \IS \times \IS \to \mathcal{B}(\mathcal{H})$ (bounded operators
on a Hilbert space) would create a ``quantum ideatic space'' in which
ideas can exist in superposition, resonance can be entangled, and
emergence can exhibit interference effects.  The spectral theory of
emergence (Section~\ref{sec:spectral-emergence}) would naturally extend
to this setting, with eigenideas becoming quantum states.

The computational implications are tantalising.  If the NP-hardness of
\textsc{MeaningNE} (Theorem~\ref{thm:ne-hard}) could be circumvented
by quantum algorithms---analogous to Shor's algorithm for factoring---then
quantum computers could in principle find optimal meaning equilibria
exponentially faster than classical computers.  Whether this is possible
depends on the structure of the quantum emergence operator, which is
entirely unexplored.
\end{game}

%% ---------------------------------------------------------------
\section{Summary of Open Problems}
\label{sec:open-summary}

Table~\ref{tab:open-problems} collects all open problems with
their status and the techniques most likely to yield progress.

\begin{table}[ht]
\centering
\small
\begin{tabular}{@{}clll@{}}
\toprule
\textbf{\#} & \textbf{Problem} & \textbf{Area} & \textbf{Likely Techniques} \\
\midrule
1  & Classification of finite $\IS$    & Algebra     & Exhaustive enumeration, GAP \\
2  & Free ideatic space                & Algebra     & Universal constructions \\
3  & Associativity conditions          & Algebra     & Operad theory \\
4  & Uniqueness of BNE                 & Game theory & Supermodularity \\
5  & ESS in meaning games              & Game theory & Replicator dynamics \\
6  & Price of anarchy                  & Game theory & Smoothness framework \\
7  & Emergence rate function           & Info.\ theory & Ergodic theory \\
8  & Source coding theorem             & Info.\ theory & Rate-distortion \\
9  & Channel coding theorem            & Info.\ theory & Random coding \\
10 & Fundamental group of~$\IS$        & Topology    & Algebraic topology \\
11 & Homology of~$\IS$                 & Topology    & Persistent homology \\
12 & Fixed-point theorems              & Topology    & Banach/Brouwer \\
13 & Resonance computation complexity  & Complexity  & Fine-grained complexity \\
14 & FPTAS for emergence               & Complexity  & Counting complexity \\
15 & Symbol grounding                  & Philosophy  & Empirical measurement \\
16 & Emergence and consciousness       & Philosophy  & IIT bridge theorems \\
17 & Quantitative cultural dynamics    & Dynamics    & ODE/PDE on $\IS$ \\
18 & Higher categories of meaning      & Category theory & $(\infty,n)$-categories \\
19 & Spectral universality             & Spectral theory & Random matrix theory \\
20 & Curvature classification          & Geometry    & Comparison geometry \\
21 & Ergodic decomposition             & Dynamics    & Ergodic theory \\
22 & Fixed-point index                 & Topology    & Lefschetz theory \\
23 & Brouwer vs.\ Hilbert in IDT      & Philosophy  & Foundations of math \\
24 & Ideatic quantum computation       & Quantum     & $C^*$-algebras \\
\bottomrule
\end{tabular}
\caption{Open problems in Idea Diffusion Theory.}
\label{tab:open-problems}
\end{table}

We close with the conviction that these twenty-four problems, taken together,
define a \emph{research programme}: a systematic attempt to
understand meaning as a mathematical phenomenon, amenable to the
same tools---algebra, topology, information theory, game theory,
computational complexity, spectral theory, ergodic theory, and the
philosophy of mathematical proof---that have illuminated the physical
sciences.

The programme has a distinctive character.  Unlike most mathematical
theories, IDT is simultaneously \emph{formal} (its theorems are
machine-verified in Lean~4) and \emph{interpretive} (its concepts
refer to human meaning, cultural practice, and social interaction).
This dual character is not a weakness but a strength: the formalism
disciplines the interpretation, and the interpretation motivates the
formalism.  The meaning games have only just begun.

\bigskip
\begin{center}
\textit{--- End of Part III ---}
\end{center}


